\RequirePackage{plautopatch}
\documentclass[uplatex, dvipdfmx, paper=a4paper, line_length=45zw, head_space=7zw, foot_space=10zw]{jlreq}
\usepackage{bxpapersize}
\usepackage[pass]{geometry}
\usepackage[T1]{fontenc}
\usepackage[deluxe, multi, jis2004]{otf}
\usepackage{IBMpsJP}
\usepackage{amsmath}
\usepackage{plext}
\usepackage{pxrubrica}
\usepackage{multicol}
\usepackage{multirow}
\usepackage{tabto}
\usepackage{gckanbun}
\usepackage{diagbox}
\usepackage{graphicx}


\setlength{\marginparpush}{0.7zw}
% \renewcommand\_{\textunderscore\allowbreak}
\DeclareEmphSequence{\usefont{T1}{lmr}{b}{n}\usefont{JY2}{hgt}{m}{n}\selectfont,\usefont{T1}{lmr}{b}{it}\usefont{JY2}{hgt}{b}{n}\selectfont}
\pretocmd{\sidenote}{\<\eghostguarded{}}{}{}
\pretocmd{\_}{\allowbreak}{}{}

\newcommand{\docDate}{Ver.1.1［2025年12月19日］}
\newcommand{\otfVersion}{2025-07-24 TeX JP org, v1.7b8 psitau, u0.32 ttk}
\newcommand{\docLicense}{\hfill{\copyright} 2025 Yarakashi Kikohshi -- \parbox[t]{66.5pt}{CC BY-NC 4.0 \\[0.5zw]\hfil\href{https://creativecommons.org/licenses/by-nc/4.0/}{\includegraphics[width=2cm]{doclicense-CC-by-nc-88x31.pdf}}\hfil}}
\newcommand{\versionWarn}{\hfil\textcolor{red}{\gtfamily\sffamily
！{\<}対象としている{japanese-otf}のバージョンは最新ではありません！}\hfil}

\makeatletter
\newcommand{\toc@dot@small}{\begin{tikzpicture} \useasboundingbox (0, 0) rectangle (4pt, 0); \fill  (2.4pt ,0.37zw) circle[radius=0.08zw]; \end{tikzpicture}}
\def\@dottedtocline#1#2#3#4#5{%
  \ifnum #1>\c@tocdepth \else
    \jlreq@set@top@contents{#1}%
    \vskip\toclineskip
    {\leftskip #2\relax \rightskip \@tocrmarg \parfillskip -\rightskip
    \parindent #2\relax\@afterindenttrue
    \interlinepenalty\@M
    \leavevmode
    \@lnumwidth #3\relax
    \@tempcnta=#1\relax
    \advance\@tempcnta by -\jlreq@top@contents
    \@tempdima=1\jlreq@mol
    \multiply \@tempdima by \@tempcnta
    \advance\leftskip \@lnumwidth \hbox{}\hskip -\leftskip
    \advance\leftskip\@tempdima
    {#4}\nobreak
    \leaders\hbox{$\m@th\mkern \@dotsep mu$\null\inhibitglue \toc@dot@small\inhibitglue\null$\m@th\mkern \@dotsep mu$}%
    % \leaders\hbox{$\m@th\mkern \@dotsep mu$\null\inhibitglue ・\inhibitglue\null$\m@th\mkern \@dotsep mu$}%
    \hfill\nobreak
    \hb@xt@\@pnumwidth{\hss\normalfont\normalcolor \tatechuyoko*{#5}}%
    \par}%
  \fi
}
\makeatother

\usepackage{jslogo}
\usepackage{hologo}
\usepackage{array}
\newcolumntype{C}[1]{>{\centering\arraybackslash}p{#1}}
\usepackage{booktabs}
\usepackage[inline]{enumitem}
\setlist[enumerate,1]{
  label=\protect\ajMaru{\arabic{*}},
  labelsep = 1em,
  leftmargin = 3em
}
\setlist[description,1]{
  font=\ttfamily,
  % font=\gtfamily\mdseries,
  labelindent=!,
  itemindent=-8zw,
  leftmargin=10zw
}
\setlist[itemize, 1]{
  label=\raisebox{0.2zw}{\textbullet},
  labelsep = 0.3zw
}
\setlist[itemize, 2]{
  % label=\raisebox{0.2zw}{\textbullet},
  labelsep = 0.3zw
}
\usepackage{bxghost}
\usepackage{fancyhdr}
\renewcommand{\headrulewidth}{0.0pt}
\usepackage{verbatim}
\usepackage{tablericons}
\usepackage{tcolorbox}
\tcbuselibrary{skins, breakable}
\usepackage{xcolor}
\usepackage{ninecolors}
\usepackage{hyperref}
\hypersetup{colorlinks,urlcolor=blue3, linkcolor=blue3}

\NewDocumentEnvironment{booktabular}{ m s }{\begin{tabular}[t]{#1}\toprule}{\bottomrule\end{tabular}}
\NewExpandableDocumentCommand{\thead}{ O{1} m }{\multicolumn{#1}{c}{\small\sffamily\gtfamily #2}}
\NewDocumentCommand{\textlrangle}{ s m }{{\rmfamily\mdseries\textlangle\nobreak\IfBooleanTF{#1}{\jpincmd{#2}}{#2}\nobreak\textrangle}}
\NewDocumentCommand{\jpincmd}{ s m }{\IfBooleanT{#1}{\inhibitkanjiskip}{\jafontsize{8.5pt}{15pt}\selectfont #2}\IfBooleanT{#1}{\inhibitkanjiskip}}
\newcommand{\inhibitkanjiskip}{\mbox{}}
\newcommand{\terminalCMD}[1]{\texttt{#1}}

\newcommand{\ptmfamily}{\fontfamily{ptm}\selectfont}
\newcommand{\ShiftJIS}{{\ptmfamily Shift\nobreak\texttt{\raisebox{0.2zw}{\textunderscore}}\nobreak JIS}}
\NewDocumentCommand{\AdobeJapan}{ o }{{\ptmfamily Adobe-Japan1\IfNoValueF{#1}{-#1}}}
\newcommand{\CTAN}{{\ptmfamily CTAN}}
\newcommand{\UCN}[1]{\texttt{U+#1}}
\newcommand{\CIDN}[1]{\texttt{CID+#1}}
\newcommand{\subCIDN}[1]{\textsubscript{(\texttt{CID+#1})}}
\newcommand{\env}[1]{\texttt{#1}環境}

\newcommand{\氏}{{\footnotesize 氏}}
\newcommand{\LuaTeX}{\hologo{LuaTeX}}
\newcommand{\TeXLive}{{\TeX} Live}
\newcommand{\TeXWiki}{{\TeX} Wiki}
\newcommand{\GitHub}{\textsf{\bfseries GitHub}}
\newcommand{\Unicode}{\textsf{Unicode}}
\newcommand{\dvipdfmx}{\terminalCMD{dvipdfmx}}
\newcommand{\DIVPDFMx}{DVIPDFM\nobreak\( x \)}
\NewDocumentCommand{\upLaTeXdvipdfmx}{ s }{\IfBooleanTF{#1}{\upLaTeX*\inhibitkanjiskip＋\inhibitkanjiskip\terminalCMD{dvipdfmx}}{\upLaTeX\inhibitkanjiskip＋\inhibitkanjiskip\terminalCMD{dvipdfmx}}}

\newcommand{\filePath}[1]{\eghostguarded{\texttt{#1}}}
\newcommand{\inlinecode}[1]{\eghostguarded{\fboxsep=1pt\fcolorbox{gray8}{white}{\texttt{#1}}}}
\newcommand{\kanjiVS}[2]{\inlinecode{\jpincmd{#1}\textbackslash symbol\{"#2\}}}
\newcommand{\kanjiUTF}[2]{\inlinecode{\jpincmd{#1}\textlrangle{\UCN{#2}}}}
\newcommand{\subUCN}[1]{\nobreak\textsubscript{(\texttt{U+#1})}}
\newcommand{\vspaceLine}{\vspace{-0.5zw}}

\newcommand{\jpinurl}{\jafontsize{8.5pt}{15pt}\selectfont}
\newcommand{\itemSet}[2]{\(\left. \parbox{#1}{#2} \right\}\)}
\newcommand{\texqa}[1]{\sidenote{\href{https://okumuralab.org/tex/mod/forum/discuss.php?d=#1}{{\TeX} Forum\,\##1}}}
\newcommand{\texqawithoutSN}[1]{\href{https://okumuralab.org/tex/mod/forum/discuss.php?d=#1}{{\TeX} Forum\,\##1}}
\newcommand{\otfUpdate}[2]{\sidenote{\parbox[t]{10zw}{\href{https://okumuralab.org/tex/mod/forum/discuss.php?d=#2}{{\TeX} Forum\, \#\textsf{#2}}\\ (#1)}}}
\newcommand{\textfrac}[2]{\( \bigl(\displaystyle\frac{\raisebox{-0.3zw}{\text{\scriptsize\gtfamily #1}}}{\raisebox{0.3zw}{\text{\scriptsize\gtfamily #2}}}\bigr) \)}
\NewDocumentCommand{\middots}{ O{1zw} }{\raisebox{#1}{\dots}}
\newcommand{\brtab}[1]{\par\tabto{#1}}

\makeatletter
\let\@@upLaTeX\upLaTeX
\RenewDocumentCommand{\upLaTeX}{ s }{\IfBooleanTF{#1}{({\kern-.1em}u{\kern-.1em}){\kern-.03em}\pLaTeX}{\@@upLaTeX}}
\makeatother

\ExplSyntaxOn
\newlength{\urltabprewidth}

\NewDocumentCommand{\urltab}{ o m }{
  \\\icon{arrow-forward}\IfNoValueF{#1}{\icon{#1}}
  \settowidth{\urltabprewidth}{\icon{arrow-forward}\IfNoValueF{#1}{\icon{#1}}}
  \parbox[t]{\dimexpr \linewidth - \urltabprewidth}
    {\setlength{\baselineskip}{1zw}#2}
}

\newenvironment{texsource}
  {\endgraf\noindent\small\verbatim}
  {\endverbatim}

\newenvironment{terminal}
  {\endgraf\noindent\small\verbatim}
  {\endverbatim}

\NewTColorBox{demosource}{ s o }{colback=gray9, left=0.5zw, top=0pt, bottom=0pt, after~skip=0.7zw, before~skip=0.7zw, sidebyside, fontupper=\ttfamily\small, fontlower=\rmfamily\normalsize, IfBooleanTF={#1}{halign~lower=center, right*=0.7zw, sidebyside~gap=0.7zw}{sidebyside~gap=1zw}, IfNoValueTF={#2}{righthand~ratio=0.4,}{righthand~ratio=#2}, colbacklower=white, skin=bicolor, colframe=gray9, leftlower=0.2zw, }
\newtcolorbox{tcb@texsource}{enhanced~jigsaw, frame~empty, colback=gray9, left=1zw, top=0pt, bottom=0pt, after~skip=0.7zw, before~skip=0.7zw, breakable=true}
\newtcolorbox{tcb@terminal}{enhanced~jigsaw, frame~empty, colback=gray2, coltext=white, left=0.5zw, top=0pt, bottom=0pt, after~skip=0.7zw, before~skip=0.7zw, breakable=true}
\newtcolorbox{directory}{before~upper=\setlength{\baselineskip}{13pt}, frame~empty, colback=green2, coltext=white, left=0.5zw, top=0pt, bottom=0pt, after~skip=0.7zw, before~skip=0.7zw, fontupper=\ttfamily}

\AddToHook{env/texsource/before}{\begin{tcb@texsource}}
\AddToHook{env/texsource/after}{\end{tcb@texsource}}
\AddToHook{env/terminal/before}{\begin{tcb@terminal}}
\AddToHook{env/terminal/after}{\end{tcb@terminal}}

% \AddToHook{env/center/before}{\medskip}
\AddToHook{env/center/after}{\bigskip}

\NewDocumentCommand{ \cmd }{ s m }{
  \IfBooleanTF{#1}
    {\uline{\eghostguarded{\texttt{\textbackslash #2}}}}
    {\eghostguarded{\texttt{\textbackslash #2}}}
}
\NewDocumentCommand{ \pkg }{ s t{+} m }{
  \IfBooleanTF{#1}
    {\textsf{#3}}
    {
      \IfBooleanTF{#2}
        {\href{https:\/\/ctan.org\/pkg\/#3}{\textsf{#3}}}
        {\textsf{#3}}
      パッケージ
    }
}
\NewDocumentCommand{ \cls }{ s t{+} m }{
  \IfBooleanTF{#1}
    {\textsf{#3}}
    {
      \IfBooleanTF{#2}
        {\href{https:\/\/ctan.org\/pkg\/#3}{\textsf{#3}}}
        {\textsf{#3}}
      文書クラス
    }
}
\NewDocumentCommand{ \opt }{ s m }{
  \IfBooleanTF{#1}
    {\texttt{#2}}
    {
      \texttt{#2}%
      オプション
    }
}
\NewDocumentCommand{\icon}{ m O{1.2zw} }{
  \eghostguarded{}
  \tablericon[height={#2/-0.3zw}]{
    \str_case:nnF { #1 }
    {
      { pdf }    { file -type -pdf }
      { text }   { file -text }
      % { text }   { file -typography }
      { github } { brand -github }
      { link }   { external -link }
    }{ #1 }
  }
  \nobreakspace
}
\ExplSyntaxOff

\NewPageStyle{cover}{nombre={}}


\ExplSyntaxOn

\NewDocumentCommand{ \ajFS }{ s m m }{%
  \IfBooleanT{#1}%
    {\colorbox{gray9}}%
  {\ibmfamily\ajFrac{#2}{#3}, \ajFrac*{#2}{#3}}%
  }%

\NewDocumentCommand{ \ajAS }{ s m }
{%
  {\ibmfamily \ajArrow{#2}}%
  &%
  \texttt{#2}%
  \IfBooleanF{#1}%
    {\hfill{\scriptsize(\cmd{ajArrow#2})}}%
}

\NewDocumentCommand{\ajLigKumimoji}{ s m m }
{%
  #2%
  & {\ibmfamily\ajLig{#2}}%
  \IfBooleanT{#1}{& {\ibmfamily\ajLig{#2*}}}%
  & {\ibmfamily\raisebox{0.1zw}{\parbox<t>{1zw}{\ajLig{#2}}}}%
  \IfBooleanT{#1}{& {\ibmfamily\raisebox{0.1zw}{\parbox<t>{1zw}{\ajLig{#2*}}}}}%
}%
\newcommand{\ajLigPunc}[1]{#1 & {\ibmfamily\ajLig{#1}}}
\NewDocumentCommand{ \ajLigKana }{ s m }{#2 & \IfBooleanT{#1}{\ibmfamily}\ajLig{#2}}
\NewDocumentCommand{ \ajPS }{ s t{+} m }
{
  {\ibmfamily\ajPICT{#3}}
  &
  \IfBooleanTF{#2}
    {
      \texttt{#3}
      \IfBooleanF{#1}
        {\hfill{\scriptsize(\cmd{ajPICT{\jpincmd* {#3}}})}}
    }
    {
      \texttt{#3}
      \IfBooleanF{#1}
        {\hfill{\scriptsize(\cmd{ajPICT#3})}}
    }
}
\ExplSyntaxOff

\newcommand{\tabularColor}{gray9}
\renewcommand{\thepage}{\large\ibmfamily\ajLabel\ajMaru{page}}

\listfiles

\usepackage{ajmacros-cmd}
\newcommand{\phrase}[1]{\texttt{#1}}
\newcommand{\optast}{\eghostguarded{\texttt{\textasteriskcentered}}}
\ExplSyntaxOn
\NewDocumentCommand{ \ajTblCntrNum }{ m s m }
{
  \IfBooleanTF{ #2 }
    { \texttt{\textbackslash #1*} }
    { \texttt{\textbackslash #1} }
  &
  \IfBooleanTF{ #2 }
    { \ibmfamily \ajCounterSample{#1}*{#3} }
    { \ibmfamily \ajCounterSample{#1}{#3} }
}
\ExplSyntaxOff
\ajCIDVarDef{高}{8705}
\ajUTFVarDef{漢}{fa47}
\ajCIDVarList{
  吉=13706
  教=8471
  鴎=7646
  崎=14290
  第=13910角=13682
  {ばく}=15192
  {Saki2}=17009
  {間2}=13693
  {nabe1}=6929
  {nabe2}=14235
  {nabe3}=14236
  {nabe4}=14237
  {nabe5}=14238
  {nabe6}=14239
  {nabe7}=14240
  {nabe8}=20234
  }
\ajUTFVarList{間=9592
  黒=9ed1
  黄=9EC3
  内=5167
  ♪=266c
  }
\title{\pkg{ajmacros}非公式マニュアル}
\author{やらかしの貴公子}
\date{\docDate}
\begin{document}
\savegeometry{body}
% \newgeometry{textwidth=45zw, textheight=600pt, marginparwidth=0pt}

\thispagestyle{fancy}
\fancyhead{}
\fancyfoot{\docLicense}
\pagestyle{cover}

\maketitle

\begin{abstract}
  これはユーザー向けの\pkg{ajmacros}の非公式マニュアルです。

  \vspace{0.7zw}
  \noindent{\centering
  \textgt{\small 対象バージョン}：%
  2025-07-24 TeX JP org, v1.7b8 psitau, u0.32 ttk%
  （コミュニティ版）}
\end{abstract}

{
\begin{multicols*}{2}
  \small
  \tableofcontents
\end{multicols*}
}

\clearpage
\loadgeometry{body}
\setcounter{page}{1}
\pagestyle{plain}

\section{\texorpdfstring{\pkg{ajmacros}}{ajmacrosパッケージ}とは}\addtocontents{section}{ajmacrosパッケージとは}

\pkg{ajmacros}は{\AdobeJapan}で規定されている約2.3万の文字を簡便に扱うことが出来るようになるパッケージです。

{\AdobeJapan}には仮名、漢字の他にもさまざまな文字・記号が収録されています。これらの文字・記号は{\ShiftJIS}では入力できない文字や{{\Unicode}}による入力であっても入力が難しい文字、{{\Unicode}}にも収録されていない文字が含まれています。

これらの文字に関して、\pkg{otf}の\cmd{CID}から入力することは可能ですが、約2.3万グリフの中から目的の文字を探し出したり数字を覚えているのは非常に難儀です。\pkg{ajmacros}はこのような煩雑さから解消します。

% \pkg{ajmacros}を読むには、文字コードがISO-2022-JPであることに注意してください。

\vfil
\subsection{パッケージ読み込み}

プリアンブルで以下のようにすることでパッケージを読み込みます。パッケージオプションはありません。

\begin{texsource}
\usepackage{ajmacros}
\end{texsource}

ただし、\pkg{ajmacros}が読み込まれるとき\pkg{otf}が読み込まれるため、\pkg{otf}だけの読み込みに終始した方が分かりやすいと思います。

\vfil
\subsection{注意}

\pkg{ajmacros}ではさまざまな記号を出力するためのコマンドがありますが、引数が必要なコマンドに与える引数が間違っている場合であってもエラーや警告が発行されません。
ただ何も無いものが出力されるだけです。

例えば、\cmd{ajLig\{\jpincmd{◇ほか}\}}（{\ibmfamily\ajLig{◇ほか}}）はありますが、\cmd{ajLig\{\jpincmd{ほか}\}}（{\ibmfamily\scalebox{0.5}{\vbox{\hbox to 2zw{ほ　}\vspace{-1zw}\hbox to 2zw{　か}}}}に相当するもの）はありません。

\begin{demosource}
  \cmd{ajLig}\{\jpincmd{◇ほか}\}はあるが、
  \cmd{ajLig}\{\jpincmd{ほか}\}はない。
  \tcblower
  {\ibmfamily\ajLig{◇ほか}}はあるが、
  \ajLig{ほか}はない。
\end{demosource}

そのため、意図せず間違った引数を与えていた場合に、その間違いに気づかない可能性があることに注意が必要です。

\vfil
\subsection{本ドキュメントで使用するフォント}

本ドキュメントの本文に利用しているフォントは原ノ味フォントです。しかし、グリフを表示するためには足りないグリフが多いため、多くにIBM Plex Sans JPを利用しています。

特にどちらのフォントを使用しているかを明記していませんが、明朝であれば原ノ味、ゴシックであればIBM Plex Sans JPであると思ってください。

\newpage

\section{漢字}

4つの漢字のコマンドが定義されています。

\begin{multicols}{2}
  \begin{itemize}
    \item \cmd{ajHashigoTaka}
      \tabto{10zw} \ajHashigoTaka (\UCN{9AD9})
    \item \cmd{ajTsuchiYoshi}
      \tabto{10zw} \ajTsuchiYoshi (\UCN{20BB7})
    \item \cmd{ajTatsuSaki}
      \tabto{10zw} \ajTatsuSaki (\UCN{FA11})
    \item \cmd{ajMayuHama}
      \tabto{10zw} \ajMayuHama (\UCN{6FF5})
  \end{itemize}
\end{multicols}

これらの文字は\cmd{CID}を利用して定義されています。そのため、同じように独自に定義してみても良いでしょう。

\subsection{文字の簡易入力定義：\texorpdfstring{\cmd{ajVar}}{\textbackslash ajVar}}

{\ShiftJIS}で直接入力できなかったりして、任意の文字をその都度{\Unicode}・CID値から指定するのは面倒です。
そこで、これらを\cmd{ajVar\{\textlrangle{文字名}\}}の形で入力できると簡単です。

これを定義するコマンドがあります。これらのコマンドはプリアンブルで使用します。

\begin{itemize}
  \item 1文字定義：
  \begin{itemize}
    \item \cmd{ajUTFVarDef\{\textlrangle{文字名}\}\{\textlrangle{{\Unicode}値}\}}
    \item \cmd{ajCIDVarDef\{\textlrangle{文字名}\}\{\textlrangle{CID値}\}}
  \end{itemize}
  \item リスト形式で複数定義：
  \begin{itemize}
    \item \cmd{ajUTFVarList\{\textlrangle{文字名1}=\textlrangle{{\Unicode}値1} \{\textlrangle{文字名2}\}=\textlrangle{{\Unicode}値2}\dots\}}
      \item \cmd{ajCIDVarList\{\textlrangle{文字名1}=\textlrangle{CID値1} \{\textlrangle{文字名2}\}=\textlrangle{CID値2}\dots\}}
  \end{itemize}
\end{itemize}

例えば、プリアンブルで以下のように定義します。\eghostguarded{\textlrangle{文字名}}が2文字以上の場合は波括弧でくくる必要があります。
また、\cmd{aj(UTF|CID)VarList}はリスト形式を採りますが、key-valueリストのように区切りにカンマ(\texttt{,})を必要としないことに注意してください。

\begin{texsource}
\ajCIDVarDef{鴎}{7646}
\ajUTFVarDef{漢}{fa47}
\ajCIDVarList{吉=13706崎=14290
  {ばく}=15192 {Saki2}=17009
  }
\ajUTFVarList{間=9592
  黒=9ed1黄=9EC3
  ♪=266c}
\end{texsource}

上記プリアンブルでは10の文字について定義しており、ドキュメント内ではこれをもとに\cmd{ajVar}を利用して文字を指定します。

\begin{demosource}
  \cmd{ajVar\{鴎\}}、\cmd{ajVar\{漢\}}。\\
  \cmd{ajVar\{吉\}}、\cmd{ajVar\{崎\}}、\cmd{ajVar\{ばく\}}、\cmd{ajVar\{Saki2\}}。\\
  \cmd{ajVar\{間\}}、\cmd{ajVar\{黒\}}、\cmd{ajVar\{黄\}}、\cmd{ajVar\{♪\}}。
  \tcblower
  \ajVar{鴎}、
  \ajVar{漢}。\\
  \ajVar{吉}、
  \ajVar{崎}、
  \ajVar{ばく}、
  \ajVar{Saki2}。\\
  \ajVar{間}、
  \ajVar{黒}、
  \ajVar{黄}、
  \ajVar{♪}。
\end{demosource}

基本的に{\AdobeJapan}が包括するすべての文字に関して構成できるため、漢字に限らないことにも留意してください。

\section{半角仮名：\texorpdfstring{\cmd{aj{\inhibitkanjiskip}半角}}{\textbackslash aj半角}}

\cmd{aj\jpincmd*{半角}}はひらがなやカタカナを半角幅にします。濁点、半濁点、小書きの仮名も範疇に入ります。一方で、漢字は半角幅にはなりません。

以下の例を見てください。

\begin{demosource}[0.3]
  {\textbackslash}aj\jpincmd*{半角}\{\jpincmd{沖縄}\}%
  で食べた%
  {\textbackslash}aj\jpincmd*{半角}\{\jpincmd{おいしい}\}%
  {\textbackslash}aj\jpincmd*{半角}\{\jpincmd{パイナップル}\}%
  と\\%
  {\textbackslash}aj\jpincmd*{半角}\{\jpincmd{すっぱい}\}%
  {\textbackslash}aj\jpincmd*{半角}\{\jpincmd{シークヮーサー}\}%
  の%
  {\textbackslash}aj\jpincmd*{半角}\{\jpincmd{ジュース}\}%
  \tcblower
  {\ibmfamily \aj半角{沖縄}で食べた\aj半角{おいしい}\aj半角{パイナップル}と%
  \aj半角{すっぱい}\aj半角{シークヮーサー}の\aj半角{ジュース}}
\end{demosource}

これらの半角幅になっている文字は、単に幅を半角幅に変形しているのではなく、半角幅のグリフを利用していることに留意してください。

\section{数値表記変換コマンド}

数値表記変換コマンドは\cmd{aj\textlrangle{囲み種}\textlrangle{文字種}\{\textlrangle{値}\}}の形式で囲み種と文字種を指定します。
代表的な囲み種と文字種はそれぞれ以下のようになっています。

\medskip
\hfil
\begin{minipage}[t]{0.45\linewidth}
  \textlrangle{囲み種}：\vspace{-0.8zw}
  \begin{itemize}
    \item 丸囲み：{\ibmfamily \ajMaru{1}}\tabto{7zw}\phrase{Maru}
    \item 四角：{\ibmfamily \ajKaku{1}}\tabto{7zw}\phrase{Kaku}
    \item 丸四角：{\ibmfamily \ajMaruKaku{1}}\tabto{7zw}\phrase{MaruKaku}
    \item 括弧囲み：{\ibmfamily \ajKakko{1}}\tabto{7zw}\phrase{Kakko}
    \item 白抜き
      \begin{itemize}
        \item 白抜き丸囲み：{\ibmfamily \ajKuroMaru{1}}\tabto{9zw}\phrase{KuroMaru}
        \item 白抜き四角：{\ibmfamily \ajKuroKaku{1}}\tabto{9zw}\phrase{KuroKaku}
        \item 白抜き丸四角：{\ibmfamily \ajKuroMaruKaku{1}}\tabto{9zw}\phrase{KuroMaruKaku}
      \end{itemize}
  \end{itemize}
\end{minipage}
\hfil
\begin{minipage}[t]{0.3\linewidth}
  \textlrangle{文字種}：\vspace{-0.8zw}
  \begin{itemize}
    \item アラビア数字 \tabto{9zw}-
    \item ラテン文字小文字 \tabto{9zw}\phrase{alph}
    \item ラテン文字大文字 \tabto{9zw}\phrase{Alph}
    \item ひらがな \tabto{9zw}\phrase{Hira}
    \item カタカナ \tabto{9zw}\phrase{Kata}
    \item 曜日 \tabto{9zw}\phrase{Yobi}
  \end{itemize}
\end{minipage}
\hfil
\bigskip

以上から、42 (7\textsubscript{\textlrangle{囲み種}}\texttimes 6\textsubscript{\textlrangle{文字種}})のコマンドが利用できます。例えば丸囲みアラビア数字は\cmd{ajMaru}、四角囲みラテン文字は\cmd{ajKakualph}と表現されます。また、他のコマンドに関しても以下のように組み合わせることで表現できます。

\begin{demosource}
  \cmd{ajMaru\{57\}}、
  \cmd{ajKakualph\{7\}}、
  \cmd{ajMaruKakuAlph\{17\}}、\\
  \cmd{ajKakkoHira\{26\}}、
  \cmd{ajKuroMaruKata\{42\}}、\\
  \cmd{ajKuroKakuYobi\{1\}}、
  \cmd{ajKuroMaruKaku*\{3\}}
  \tcblower
  \ibmfamily
  \ajMaru{57}、
  \ajKakualph{7}、
  \ajMaruKakuAlph{17}、
  \ajKakkoHira{26}、
  \ajKuroMaruKata{42}、
  \ajKuroKakuYobi{1}、
  \ajKuroMaruKaku*{3}
\end{demosource}

これらの数値表記変換コマンドは、CID値が連続することを利用して、下限となるグリフのCID値からのズレを計算することでグリフを指定します。
このとき、指定するズレの量には上限がありません。
そのため、最大値を超えるような数字が指定された場合、まったく意図しない文字が出力されるようになります。
\warichu{ただし、どのグリフもCID値が連続しているわけではないため、適宜、条件分岐によって値が飛ぶようになっています。}

一方で、類似グリフのCID値が連続するような場合では少し便利です。例えば、丸囲みラテン文字は小文字から大文字が連続したCID値のため、\cmd{ajMarualph}で26を超えたときでも丸囲みラテン文字大文字になります。

\begin{demosource}
  {\textbackslash}ajMarualph\{24\},
  {\textbackslash}ajMarualph\{25\},
  {\textbackslash}ajMarualph\{26\},\\
  {\textbackslash}ajMarualph\{27\},
  {\textbackslash}ajMarualph\{28\},
  {\textbackslash}ajMarualph\{29\}.
  \tcblower
  \ibmfamily
  \ajMarualph{24},
  \ajMarualph{25},
  \ajMarualph{26},
  \ajMarualph{27},
  \ajMarualph{28},
  \ajMarualph{29}.
\end{demosource}

\newpage
\subsection{アラビア数字}

それぞれに{\optast}付きのコマンド変種があります。これらは1桁のアラビア数字の先頭に0が加えられたグリフになっています。

\begin{center}
  \begin{booktabular}{p{12zw}lcc}
    \thead{コマンド}             & \thead{例}           & \thead{最小値} & \thead{最大値}    \\
    \midrule
    \ajTblCntrNum{ajMaru}          {0,1,2-50,51,52-99,100} & 0              & 100            \\
    \ajTblCntrNum{ajMaru}*         {0,1,2-50,51,52-99,100} & 0              & 100            \\
    \ajTblCntrNum{ajKuroMaru}      {0,1,2-50,51,52-99,100} & 0              & 100            \\
    \ajTblCntrNum{ajKuroMaru}*     {0,1,2-50,51,52-99,100} & 0              & 100            \\
    \ajTblCntrNum{ajKaku}          {0,1,2-50,51,52-99,100} & 0              & 100            \\
    \ajTblCntrNum{ajKaku}*         {0,1,2-50,51,52-99,100} & 0              & 100            \\
    \ajTblCntrNum{ajKuroKaku}      {0,1,2-50,51,52-99,100} & 0              & 100            \\
    \ajTblCntrNum{ajKuroKaku}*     {0,1,2-50,51,52-99,100} & 0              & 100            \\
    \ajTblCntrNum{ajMaruKaku}      {0,1,2-50,51,52-99,100} & 0              & 100            \\
    \ajTblCntrNum{ajMaruKaku}*     {0,1,2-50,51,52-99,100} & 0              & 100            \\
    \ajTblCntrNum{ajKuroMaruKaku}  {0,1,2-50,51,52-99,100} & 0              & 100            \\
    \ajTblCntrNum{ajKuroMaruKaku}* {0,1,2-50,51,52-99,100} & 0              & 100            \\
    \ajTblCntrNum{ajKakko}         {0,1,2-50,51,52-99,100} & 0              & 100            \\
    \ajTblCntrNum{ajKakko}*        {0,1,2-50,51,52-99,100} & 0              & 100            \\
  \end{booktabular}
\end{center}

\vfil
\subsection{ラテン文字}

\begin{center}
  \begin{booktabular}{p{12zw}lcc}
    \thead{コマンド}             & \thead{例}           & \thead{最小値} & \thead{最大値} \\
    \midrule
    \ajTblCntrNum{ajMarualph}         {1,2,3-15,16,17-25,26} & 1          & 26            \\
    \ajTblCntrNum{ajMaruAlph}         {1,2,3-15,16,17-25,26} & 1          & 26            \\
    \ajTblCntrNum{ajKuroMarualph}     {1,2,3-15,16,17-25,26} & 1          & 26            \\
    \ajTblCntrNum{ajKuroMaruAlph}     {1,2,3-15,16,17-25,26} & 1          & 26            \\
    \ajTblCntrNum{ajKakualph}         {1,2,3-15,16,17-25,26} & 1          & 26            \\
    \ajTblCntrNum{ajKakuAlph}         {1,2,3-15,16,17-25,26} & 1          & 26            \\
    \ajTblCntrNum{ajKuroKakualph}     {1,2,3-15,16,17-25,26} & 1          & 26            \\
    \ajTblCntrNum{ajKuroKakuAlph}     {1,2,3-15,16,17-25,26} & 1          & 26            \\
    \ajTblCntrNum{ajMaruKakualph}     {1,2,3-15,16,17-25,26} & 1          & 26            \\
    \ajTblCntrNum{ajMaruKakuAlph}     {1,2,3-15,16,17-25,26} & 1          & 26            \\
    \ajTblCntrNum{ajKuroMaruKakualph} {1,2,3-15,16,17-25,26} & 1          & 26            \\
    \ajTblCntrNum{ajKuroMaruKakuAlph} {1,2,3-15,16,17-25,26} & 1          & 26            \\
    \ajTblCntrNum{ajKakkoalph}        {1,2,3-15,16,17-25,26} & 1          & 26            \\
    \ajTblCntrNum{ajKakkoAlph}        {1,2,3-15,16,17-25,26} & 1          & 26            \\
  \end{booktabular}
\end{center}

\newpage
\subsection{ひらがな・カタカナ}

\begin{center}
  \begin{booktabular}{p{12zw}lcc}
    \thead{コマンド}             & \thead{例}           & \thead{最小値} & \thead{最大値} \\
    \midrule
    \ajTblCntrNum{ajMaruHira}         {1,2,3-25,26,27-47,48} & 1          & 48            \\
    \ajTblCntrNum{ajMaruKata}         {1,2,3-25,26,27-47,48} & 1          & 48            \\
    \ajTblCntrNum{ajKuroMaruHira}     {1,2,3-25,26,27-47,48} & 1          & 48            \\
    \ajTblCntrNum{ajKuroMaruKata}     {1,2,3-25,26,27-47,48} & 1          & 48            \\
    \ajTblCntrNum{ajKakuHira}         {1,2,3-25,26,27-47,48} & 1          & 48            \\
    \ajTblCntrNum{ajKakuKata}         {1,2,3-25,26,27-47,48} & 1          & 48            \\
    \ajTblCntrNum{ajKuroKakuHira}     {1,2,3-25,26,27-47,48} & 1          & 48            \\
    \ajTblCntrNum{ajKuroKakuKata}     {1,2,3-25,26,27-47,48} & 1          & 48            \\
    \ajTblCntrNum{ajMaruKakuHira}     {1,2,3-25,26,27-47,48} & 1          & 48            \\
    \ajTblCntrNum{ajMaruKakuKata}     {1,2,3-25,26,27-47,48} & 1          & 48            \\
    \ajTblCntrNum{ajKuroMaruKakuHira} {1,2,3-25,26,27-47,48} & 1          & 48            \\
    \ajTblCntrNum{ajKuroMaruKakuKata} {1,2,3-25,26,27-47,48} & 1          & 48            \\
    \ajTblCntrNum{ajKakkoHira}        {1,2,3-25,26,27-47,48} & 1          & 48            \\
    \ajTblCntrNum{ajKakkoKata}        {1,2,3-25,26,27-47,48} & 1          & 48            \\
  \end{booktabular}
\end{center}

濁点、半濁点、小書きの囲み文字はありません。以下に通しのグリフを掲載します。

\begin{description}[rightmargin=3zw]
  \item[ひらがな：] {\ibmfamily \ajMaruHira{1}, \ajMaruHira{2}, \ajMaruHira{3}, \ajMaruHira{4}, \ajMaruHira{5}, \ajMaruHira{6}, \ajMaruHira{7}, \ajMaruHira{8}, \ajMaruHira{9}, \ajMaruHira{10}, \ajMaruHira{11}, \ajMaruHira{12}, \ajMaruHira{13}, \ajMaruHira{14}, \ajMaruHira{15}, \ajMaruHira{16}, \ajMaruHira{17}, \ajMaruHira{18}, \ajMaruHira{19}, \ajMaruHira{20}, \ajMaruHira{21}, \ajMaruHira{22}, \ajMaruHira{23}, \ajMaruHira{24}, \ajMaruHira{25}, \ajMaruHira{26}, \ajMaruHira{27}, \ajMaruHira{28}, \ajMaruHira{29}, \ajMaruHira{30}, \ajMaruHira{31}, \ajMaruHira{32}, \ajMaruHira{33}, \ajMaruHira{34}, \ajMaruHira{35}, \ajMaruHira{36}, \ajMaruHira{37}, \ajMaruHira{38}, \ajMaruHira{39}, \ajMaruHira{40}, \ajMaruHira{41}, \ajMaruHira{42}, \ajMaruHira{43}, \ajMaruHira{44}, \ajMaruHira{45}, \ajMaruHira{46}, \ajMaruHira{47}, \ajMaruHira{48}.}
  \item[カタカナ：] {\ibmfamily \ajMaruKata{1}, \ajMaruKata{2}, \ajMaruKata{3}, \ajMaruKata{4}, \ajMaruKata{5}, \ajMaruKata{6}, \ajMaruKata{7}, \ajMaruKata{8}, \ajMaruKata{9}, \ajMaruKata{10}, \ajMaruKata{11}, \ajMaruKata{12}, \ajMaruKata{13}, \ajMaruKata{14}, \ajMaruKata{15}, \ajMaruKata{16}, \ajMaruKata{17}, \ajMaruKata{18}, \ajMaruKata{19}, \ajMaruKata{20}, \ajMaruKata{21}, \ajMaruKata{22}, \ajMaruKata{23}, \ajMaruKata{24}, \ajMaruKata{25}, \ajMaruKata{26}, \ajMaruKata{27}, \ajMaruKata{28}, \ajMaruKata{29}, \ajMaruKata{30}, \ajMaruKata{31}, \ajMaruKata{32}, \ajMaruKata{33}, \ajMaruKata{34}, \ajMaruKata{35}, \ajMaruKata{36}, \ajMaruKata{37}, \ajMaruKata{38}, \ajMaruKata{39}, \ajMaruKata{40}, \ajMaruKata{41}, \ajMaruKata{42}, \ajMaruKata{43}, \ajMaruKata{44}, \ajMaruKata{45}, \ajMaruKata{46}, \ajMaruKata{47}, \ajMaruKata{48}.}
\end{description}

\vfil
\subsection{曜日}

\begin{center}
  \begin{booktabular}{p{12zw}lccc}
    \thead{コマンド}             & \thead{例}         & \thead{最小値} & \thead{最大値} \\
    \midrule
    \ajTblCntrNum{ajMaruYobi}         {1,2,3,4,5,6,7} & 1          & 7             \\
    \ajTblCntrNum{ajKuroMaruYobi}     {1,2,3,4,5,6,7} & 1          & 7             \\
    \ajTblCntrNum{ajKakuYobi}         {1,2,3,4,5,6,7} & 1          & 7             \\
    \ajTblCntrNum{ajKuroKakuYobi}     {1,2,3,4,5,6,7} & 1          & 7             \\
    \ajTblCntrNum{ajMaruKakuYobi}     {1,2,3,4,5,6,7} & 1          & 7             \\
    \ajTblCntrNum{ajKuroMaruKakuYobi} {1,2,3,4,5,6,7} & 1          & 7             \\
    \ajTblCntrNum{ajKakkoYobi}        {1,2,3,4,5,6,7} & 1          & 7             \\
  \end{booktabular}
\end{center}

\newpage
\subsection{その他}

\cmd{ajApostrophe}と\cmd{ajYear}は同じです。
\cmd{ajRoman*}は4のみグリフが異なります。

\cmd{ajSquareMark}には\cmd{ajSquareMark+}という変種もあります。
これは、偶数を\cmd{ajSquareMark}に、奇数を\cmd{ajSquareMark*}にします。
\cmd{ajSquareMark}とその変種の{\optast}付きが出力する記号に関して、6番目の記号が対応していないことに注意してください。

\begin{center}
  \begin{booktabular}{p{11zw}lcc}
    \thead{コマンド}             & \thead{例}           & \thead{最小値} & \thead{最大値} \\
    \midrule
    \ajTblCntrNum{ajNijuMaru}      {1,2,3,4,5,6,7,8,9,10}      & 1   & 10  \\
    \ajTblCntrNum{ajPeriod}        {1,2,3,4,5,6,7,8,9,10,11}   & 1   & 11  \\
    \ajTblCntrNum{ajSlanted}       {0,1,2-29,30,31-58,59}      & 0   & 59  \\
    \ajTblCntrNum{ajSlanted}*      {0,1,2-29,30,31-58,59}      & 0   & 59  \\
    \ajTblCntrNum{ajCross}         {1,2,3,4,5,6,7,8,9}         & 1   & 9   \\
    \ajTblCntrNum{ajApostrophe}    {0,1,2-50,51,52-98,99}      & 0   & 99  \\
    \ajTblCntrNum{ajYear}          {0,1,2-50,51,52-98,99}      & 0   & 99  \\
    \ajTblCntrNum{ajHasenKakuAlph} {1,2,3,4,5,6}               & 1   & 6   \\
    \ajTblCntrNum{ajMaruKansuji}   {1,2,3,4,5,6,7,8,9,10}      & 1   & 10  \\
    \ajTblCntrNum{ajKakkoKansuji}  {1,2,3,4,5,6-18,19,20}      & 1   & 20  \\
    \ajTblCntrNum{ajroman}         {1,2,3,4,5-12,13,14,15}     & 1   & 15  \\
    \ajTblCntrNum{ajRoman}         {1,2,3,4,5-12,13,14,15}     & 1   & 15  \\
    \ajTblCntrNum{ajRoman}*        {1,2,3,4,5-12,13,14,15}     & 1   & 15  \\
    % \ajTblCntrNum{ajKakkoroman}    {1,2,3,4,5-12,13,14,15}     & 1   & 15  \\
    \cmd{ajKakkoroman}            & {\ibmfamily \CID{9984}\,\CID{9985}\,\CID{9986}\,\CID{9987}\,\CID{9988}\dots\CID{10000}\,\CID{10001}\,\CID{10002}\,\CID{10003}}
                                                              & 1   & 15  \\
    % \ajTblCntrNum{ajKakkoRoman}    {1,2,3,4,5-12,13,14,15}     & 1   & 15  \\
    \cmd{ajKakkoRoman}            & {\ibmfamily \CID{9974}\,\CID{9975}\,\CID{9976}\,\CID{9977}\,\CID{9978}\dots\CID{9994}\,\CID{9995}\,\CID{9996}\,\CID{9997}\,\CID{9998}}
                                                              & 1   & 15  \\
    \ajTblCntrNum{ajRecycle}       {0,1,2,3,4,5,6,7,8,9,10,11} & 0   & 11  \\
    \ajTblCntrNum{ajSquareMark}    {1,2,3,4,5,6,7,8,9,10}      & 1   & 10  \\
    \ajTblCntrNum{ajSquareMark}*   {1,2,3,4,5,6,7,8,9}         & 1   & 9   \\
    \ajTblCntrNum{ajHishi}         {1,2,3,4}                   & 1   & 4   \\
  \end{booktabular}
\end{center}

IBM Plex Sans JPの不備のため、本来\cmd{ajRoman*}の{\ibmfamily\CID{9883}}に丸括弧はないはずですが、なぜか囲まれたグリフになっています。
加えて、\cmd{ajKakkoroman}および\cmd{ajKakkoRoman}は、グリフ順序が間違っていました。そのため、適切な順序になるように調整しています。

\section{詰め数字：\texorpdfstring{\cmd{ajTsumesuji}}{\textbackslash ajTsumesuji}}

\cmd{ajTsumesuji} (\cmd{ajTumesuji})によって、詰め数字を表現できます。第1引数の値によって幅が変わります。（字幅を変形で詰めているのではなく、対応する幅のグリフがあります）

\begin{itemize}
  \item \cmd{ajTsumesuji\{1\}\{\textlrangle{digits}\}}：全角幅 \qquad {\small\gtfamily 例}： {\ibmfamily \ajTsumesuji{1}{0123456789}}
  \item \cmd{ajTsumesuji\{2\}\{\textlrangle{digits}\}}：二分幅 \qquad {\small\gtfamily 例}： {\ibmfamily \ajTsumesuji{2}{0123456789}}
  \item \cmd{ajTsumesuji\{3\}\{\textlrangle{digits}\}}：三分幅 \qquad {\small\gtfamily 例}： {\ibmfamily \ajTsumesuji{3}{0123456789}}
  \item \cmd{ajTsumesuji\{4\}\{\textlrangle{digits}\}}：四分幅 \qquad {\small\gtfamily 例}： {\ibmfamily \ajTsumesuji{4}{0123456789}}
\end{itemize}

加えて、\cmd{ajTsumesuji*}は文字数に応じて自動的に幅を変更します。
ただし、この\cmd{ajTsumesuji*\{\textlrangle{digits}\}}は5桁を超える数字は受け付けません。

また、\cmd{ajTsumesuji(*)}の第2引数は数字のみを受け付けるため、カンマ(\texttt{,})やピリオド(\texttt{.})を含めることは出来ません。

これらのコマンドは、特に縦組中に用いられる連数字（縦中横）を表現する際に利用できます。\pkg{plext}の\cmd{rensuji}や\cls{jlreq}の\cmd{tatechuyoko}を利用して、以下のようにします。

\begin{demosource}*[0.15]
  平成は{\textbackslash}rensuji\{\textbackslash{ajTsumesuji}\{4\}\{2019\}\}年
  {\textbackslash}rensuji\{\textbackslash{ajTsumesuji}\{1\}\{4\}\}月\%\\
  {\textbackslash}rensuji\{\textbackslash{ajTsumesuji}\{2\}\{30\}\}日
  に終わりを告げ、その期間は\%\\
  {\textbackslash}rensuji\{\textbackslash{ajTsumesuji}\{4\}\{11070\}\}日
  （{\textbackslash}rensuji\{\textbackslash{ajTsumesuji}*\{30\}\}年と\%\\
  {\textbackslash}rensuji\{\textbackslash{ajTsumesuji}*\{113\}\}日）となりました。
  \tcblower
  \begin{minipage}<t>[t][4zw]{15zw}
    平成は\rensuji{\ibmfamily\ajTsumesuji{4}{2019}}年\rensuji{\ibmfamily\ajTsumesuji{1}{4}}月\rensuji{\ibmfamily\ajTsumesuji{2}{30}}日に終わりを告げ、その期間は
    \rensuji{\ibmfamily\ajTsumesuji{4}{11070}}日（\rensuji{\ibmfamily\ajTsumesuji*{30}}年と\rensuji{\ibmfamily\ajTsumesuji*{113}}日）となりました。
  \end{minipage}
\end{demosource}

\section{数値表記変換コマンドや詰め数字をカウンタに利用する}

数値表記変換コマンドや詰め数字をカウンタに利用することが出来ます。これらのコマンドをカウンタに利用する場合、以下の2点について注意が必要です。

\begin{itemize}
  \item fragileなコマンドなため、\cmd{section}等の中では\cmd{protect}を前置する
  \item カウンタを引数にする場合は\cmd{ajLabel}を前置する
\end{itemize}

数値表記変換コマンドはfragileなコマンドのため、\cmd{section}内であっても\cmd{protect}を利用する必要があります。

\begin{texsource}
\section{\protect\ajMaru{5}}
\end{texsource}

また、\pkg{enumitem}を使用する場合も同様に\cmd{protect}を使用します。

\begin{texsource}
\begin{enumerate}[label=\protect\ajMaru{\arabic{*}}] ... \end{enumerate}
\end{texsource}

数値表記変換コマンドや詰め数字は引数に数値のみを採るため、カウンタを引数にする場合は\cmd{ajLabel}を前置します。これによって、\cmd{arabic\{\textlrangle{counter}\}}等と同じようにカウンタを引数に取ることが出来ます。

例えば、ページ番号を丸囲み数字にする、番号付きリストを詰め数字にするには以下のようにします。

\begin{texsource}
\renewcommand{\thepage}{\ajLabel\ajMaru{page}}
\renewcommand{\theenumi}{\ajLabel\ajTumesuji*{enumi}}
\end{texsource}

ただし、\cmd{ajLabel}が\pkg{hyperref}の更新によってはエラーが生じることがあります。

\begin{itemize}
  \item k16's note: \LaTeX で連番リストに丸数字（\cmd{ajMaru}）を使えたはずだけどhyperrefで困る話
  \urltab{\url{https://note.golden-lucky.net/2016/02/latexajmaruhyperref.html}}
\end{itemize}

\section{詰め括弧、年齢}

詰め数字を括弧で囲った詰め括弧と新聞等で見られる丸括弧に数字と「つ」や「才」等を与えた年齢の表現をします。

\begin{itemize}
  %%
  \item \cmd{ajTsumekakko\{\textlrangle{digits}\}}
  \begin{itemize}
    \item {\small\sffamily\gtfamily 横組：}
      \begin{demosource}[0.25]
        {\textbackslash}ajTsumekakko\{20\}%
        {\textbackslash}ajTsumekakko\{123\}%
        {\textbackslash}ajTsumekakko\{1998\}%
        \tcblower
        \ibmfamily
        \ajTsumekakko{20}
        \ajTsumekakko{123}
        \ajTsumekakko{1998}
      \end{demosource}
    \item {\small\sffamily\gtfamily 縦組：}
      \begin{demosource}[0.25]
        {\textbackslash}pbox<t>\{{\textbackslash}ajTsumekakko\{20\}\}\\
        {\textbackslash}pbox<t>\{{\textbackslash}ajTsumekakko\{123\}\}\\
        {\textbackslash}pbox<t>\{{\textbackslash}ajTsumekakko\{1998\}\}
        \tcblower
        \ibmfamily
        \pbox<t>{\ajTsumekakko{20}}
        \pbox<t>{\ajTsumekakko{123}}
        \pbox<t>{\ajTsumekakko{1998}}
      \end{demosource}
  \end{itemize}
  %%
  \item \cmd{ajNenrei\{\textlrangle{1, 2}\}\{\textlrangle{digits}\}}
    \begin{itemize}
      \item {\small\sffamily\gtfamily 横組：}
        \begin{demosource}[0.25]
          {\textbackslash}ajNenrei\{1\}\{3\}%
          {\textbackslash}ajNenrei\{2\}\{5\}\\%
          {\textbackslash}ajNenrei\{1\}\{13\}%
          {\textbackslash}ajNenrei\{2\}\{25\}%
          \tcblower
          \ibmfamily
          \ajNenrei{1}{3}%
          \ajNenrei{2}{5}%
          \ajNenrei{1}{13}%
          \ajNenrei{2}{25}%
        \end{demosource}
      \item {\small\sffamily\gtfamily 縦組：}
        \begin{demosource}[0.25]
          {\textbackslash}pbox<t>\{{\textbackslash}ajNenrei\{1\}\{3\}\}%
          {\textbackslash}pbox<t>\{{\textbackslash}ajNenrei\{1\}\{13\}\}、\\%
          {\textbackslash}pbox<t>\{{\textbackslash}ajNenrei\{2\}\{5\}\}%
          {\textbackslash}pbox<t>\{{\textbackslash}ajNenrei\{2\}\{25\}\}%
          \tcblower
          \ibmfamily
          \pbox<t>{\ajNenrei{1}{3}}%
          \pbox<t>{\ajNenrei{1}{13}}、%
          \pbox<t>{\ajNenrei{2}{5}}%
          \pbox<t>{\ajNenrei{2}{25}}%
        \end{demosource}
    \end{itemize}
  %%
  \item \cmd{ajNenrei\{\textlrangle{3, 4}\}\{\textlrangle{1 -- 99}\}}（縦組のみ横組では意味をなさない）
    \begin{demosource}[0.25]
      {\textbackslash}pbox<t>\{{\textbackslash}ajNenrei\{3\}\{7\}\}%
      {\textbackslash}pbox<t>\{{\textbackslash}ajNenrei\{3\}\{37\}\}、\\%
      {\textbackslash}pbox<t>\{{\textbackslash}ajNenrei\{4\}\{9\}\}%
      {\textbackslash}pbox<t>\{{\textbackslash}ajNenrei\{4\}\{59\}\}%
      \tcblower
      \ibmfamily
      \pbox<t>{\ajNenrei{3}{7}}%
      \pbox<t>{\ajNenrei{3}{37}}、%
      \pbox<t>{\ajNenrei{4}{9}}%
      \pbox<t>{\ajNenrei{4}{59}}%
    \end{demosource}
  %%
  \item \cmd{ajnenrei\{\textlrangle{1, 2, 3, 4}\}\{\textlrangle{1 -- 99}\}}（組方向に拘わらず結果は同じ）
    \begin{demosource}[0.25]
      {\textbackslash}ajnenrei\{1\}\{3\}%
      {\textbackslash}ajnenrei\{2\}\{5\}\\%
      {\textbackslash}ajnenrei\{3\}\{7\}%
      {\textbackslash}ajnenrei\{4\}\{9\}%
      \tcblower
      \ibmfamily
      \ajnenrei{1}{3}
      \ajnenrei{2}{5}
      \ajnenrei{3}{7}
      \ajnenrei{4}{9}
    \end{demosource}
\end{itemize}

\vfil
\subsection{\texorpdfstring{\cmd{ajTsumekakko}}{\textbackslash ajTsumekakko}}

\cmd{ajTsumekakko}は組方向によって結果が異なります。ただし、横組1桁の場合は結果が崩れてしまうようです。
加えて、\cmd{ajNenrei\{0\}}も同じ結果を得ます。

\begin{center}
  \begin{booktabular}{cl@{\hspace{2zw}}l}
    \thead{組方向}\ & \thead{1 -- 9} & \thead{10 --} \\
    \midrule
    \textgt{横} &
    % {\ibmfamily \ajTsumekakko{1}\,\ajTsumekakko{2}\,\ajTsumekakko{3}\middots[0.3zw]\ajTsumekakko{8}\,\ajTsumekakko{9}}&
    {\ttfamily --} &
    {\ibmfamily \ajTsumekakko{10}\,\ajTsumekakko{11}\,\ajTsumekakko{12}\middots[0.3zw]\ajTsumekakko{98}\,\ajTsumekakko{99}\,\ajTsumekakko{100}\middots[0.3zw]}
    \\
    \textgt{縦} &
    {\ibmfamily {\pbox<t>{\ajTsumekakko{1}}}\,{\pbox<t>{\ajTsumekakko{2}}}\,{\pbox<t>{\ajTsumekakko{3}}}\middots[0.5zw]{\pbox<t>{\ajTsumekakko{8}}}\,{\pbox<t>{\ajTsumekakko{9}}}} &
    {\ibmfamily \raisebox{-0.5zw}{\pbox<t>{\ajTsumekakko{10}}}\,\raisebox{-0.5zw}{\pbox<t>{\ajTsumekakko{11}}}\,\raisebox{-0.5zw}{\pbox<t>{\ajTsumekakko{12}}}\middots[0.5zw]\raisebox{-0.5zw}{\pbox<t>{\ajTsumekakko{98}}}\,\raisebox{-0.5zw}{\pbox<t>{\ajTsumekakko{99}}}\,\raisebox{-0.5zw}{\pbox<t>{\ajTsumekakko{100}}}\middots[0.5zw]\,\raisebox{-1zw}{\pbox<t>{\ajTsumekakko{1999}}}}
    \\[0.8zw]
  \end{booktabular}
\end{center}

\vfil

\subsection{\texorpdfstring{\cmd{ajNenrei}}{\textbackslash ajNenrei}（横組）}

横組の\cmd{ajNenrei\{\textlrangle{1, 2}\}}は以下のようになります。\cmd{ajNenrei\{\textlrangle{3, 4}\}}は横組では意味をなさないため掲示していません。（例：\cmd{ajNenrei\{3\}\{1\}\,→\,{\ibmfamily\ajNenrei{3}{1}}}）

\begin{center}
  \begin{booktabular}{cll}
    \diagbox[font=\sffamily\gtfamily\small]{第1引数}{第2引数} & \thead{1 -- 9} & \thead{10 -- 99} \\
    \hline
    {\texttt{1} {\scriptsize (つ)}} &
    {\ibmfamily \ajNenrei{1}{1}\,\ajNenrei{1}{2}\,\ajNenrei{1}{3}\middots[0.3zw]\ajNenrei{1}{8}\,\ajNenrei{1}{9}} &
    {\ibmfamily \ajNenrei{1}{10}\,\ajNenrei{1}{11}\,\ajNenrei{1}{12}\middots[0.3zw]\ajNenrei{1}{98}\,\ajNenrei{1}{99}}
    \\
    {\texttt{2} {\scriptsize (才)}} &
    {\ibmfamily \ajNenrei{2}{1}\,\ajNenrei{2}{2}\,\ajNenrei{2}{3}\middots[0.3zw]\ajNenrei{2}{8}\,\ajNenrei{2}{9}} &
    {\ibmfamily \ajNenrei{2}{10}\,\ajNenrei{2}{11}\,\ajNenrei{2}{12}\middots[0.3zw]\ajNenrei{2}{98}\,\ajNenrei{2}{99}}
    \\
  \end{booktabular}
\end{center}

\subsection{\texorpdfstring{\cmd{ajNenrei}}{\textbackslash ajNenrei}（縦組）}

縦組の\cmd{ajNenrei\{\textlrangle{1, 2, 3, 4}\}}は以下のようになります。\cmd{ajNenrei\{\textlrangle{3, 4}\}}の第2引数が3桁になる場合、意味をなさないため掲示していません。（例：\cmd{ajNenrei\{3\}\{100\}\,→\,{\ibmfamily\pbox<t>{\ajNenrei{3}{100}}}}）

\begin{center}
  \begin{booktabular}{ccccccc}
    \diagbox[font=\sffamily\gtfamily\small]{第1引数}{第2引数} & \thead{1 -- 9} & \thead{10} & \thead{11 -- 19} & \thead{20 -- 99} & \thead{100 --} \\
    \hline
    \raisebox{0.7zw}{\texttt{1} {\scriptsize (つ)}} &
    {\ibmfamily \pbox<t>{\ajNenrei{1}{1}}\,\pbox<t>{\ajNenrei{1}{2}}\,\pbox<t>{\ajNenrei{1}{3}}\middots\pbox<t>{\ajNenrei{1}{8}}\,\pbox<t>{\ajNenrei{1}{9}}} &
    {\ibmfamily \pbox<t>{\ajNenrei{1}{10}}} &
    {\ibmfamily \pbox<t>{\ajNenrei{1}{11}}\,\pbox<t>{\ajNenrei{1}{12}}\,\pbox<t>{\ajNenrei{1}{13}}\middots\pbox<t>{\ajNenrei{1}{18}}\,\pbox<t>{\ajNenrei{1}{19}}} &
    {\ibmfamily \pbox<t>{\ajNenrei{1}{20}}\,\pbox<t>{\ajNenrei{1}{21}}\,\pbox<t>{\ajNenrei{1}{22}}\middots\pbox<t>{\ajNenrei{1}{98}}\,\pbox<t>{\ajNenrei{1}{99}}} &
    {\ibmfamily \pbox<t>{\ajNenrei{1}{100}}\,\pbox<t>{\ajNenrei{1}{101}}\,\pbox<t>{\ajNenrei{1}{102}}\middots\pbox<t>{\ajNenrei{1}{1098}}\,\pbox<t>{\ajNenrei{1}{1099}}}
    \\
    \raisebox{0.7zw}{\texttt{2} {\scriptsize (才)}} &
    {\ibmfamily \pbox<t>{\ajNenrei{2}{1}}\,\pbox<t>{\ajNenrei{2}{2}}\,\pbox<t>{\ajNenrei{2}{3}}\middots\pbox<t>{\ajNenrei{2}{8}}\,\pbox<t>{\ajNenrei{2}{9}}} &
    {\ibmfamily \pbox<t>{\ajNenrei{2}{10}}} &
    {\ibmfamily \pbox<t>{\ajNenrei{2}{11}}\,\pbox<t>{\ajNenrei{2}{12}}\,\pbox<t>{\ajNenrei{2}{13}}\middots\pbox<t>{\ajNenrei{2}{18}}\,\pbox<t>{\ajNenrei{2}{19}}} &
    {\ibmfamily \pbox<t>{\ajNenrei{2}{20}}\,\pbox<t>{\ajNenrei{2}{21}}\,\pbox<t>{\ajNenrei{2}{22}}\middots\pbox<t>{\ajNenrei{2}{98}}\,\pbox<t>{\ajNenrei{2}{99}}} &
    {\ibmfamily \pbox<t>{\ajNenrei{2}{100}}\,\pbox<t>{\ajNenrei{2}{101}}\,\pbox<t>{\ajNenrei{2}{102}}\middots\pbox<t>{\ajNenrei{2}{1098}}\,\pbox<t>{\ajNenrei{2}{1099}}}
    \\
    \raisebox{0.7zw}{\texttt{3} {\scriptsize (つ)}} &
    {\ibmfamily \pbox<t>{\ajNenrei{3}{1}}\,\pbox<t>{\ajNenrei{3}{2}}\,\pbox<t>{\ajNenrei{3}{3}}\middots\pbox<t>{\ajNenrei{3}{8}}\,\pbox<t>{\ajNenrei{3}{9}}} &
    {\ibmfamily \pbox<t>{\ajNenrei{3}{10}}} &
    {\ibmfamily \pbox<t>{\ajNenrei{3}{11}}\,\pbox<t>{\ajNenrei{3}{12}}\,\pbox<t>{\ajNenrei{3}{13}}\middots\pbox<t>{\ajNenrei{3}{18}}\,\pbox<t>{\ajNenrei{3}{19}}} &
    {\ibmfamily \pbox<t>{\ajNenrei{3}{20}}\,\pbox<t>{\ajNenrei{3}{21}}\,\pbox<t>{\ajNenrei{3}{22}}\middots\pbox<t>{\ajNenrei{3}{98}}\,\pbox<t>{\ajNenrei{3}{99}}} &
    \texttt{--}
    \\
    \raisebox{0.7zw}{\texttt{4} {\scriptsize (才)}} &
    {\ibmfamily \pbox<t>{\ajNenrei{4}{1}}\,\pbox<t>{\ajNenrei{4}{2}}\,\pbox<t>{\ajNenrei{4}{3}}\middots\pbox<t>{\ajNenrei{4}{8}}\,\pbox<t>{\ajNenrei{4}{9}}} &
    {\ibmfamily \pbox<t>{\ajNenrei{4}{10}}} &
    {\ibmfamily \pbox<t>{\ajNenrei{4}{11}}\,\pbox<t>{\ajNenrei{4}{12}}\,\pbox<t>{\ajNenrei{4}{13}}\middots\pbox<t>{\ajNenrei{4}{18}}\,\pbox<t>{\ajNenrei{4}{19}}} &
    {\ibmfamily \pbox<t>{\ajNenrei{4}{20}}\,\pbox<t>{\ajNenrei{4}{21}}\,\pbox<t>{\ajNenrei{4}{22}}\middots\pbox<t>{\ajNenrei{4}{98}}\,\pbox<t>{\ajNenrei{4}{99}}} &
    \texttt{--}
    \\
  \end{booktabular}
\end{center}

\subsection{\texorpdfstring{\cmd{ajnenrei}}{\textbackslash ajnenrei}}

\cmd{ajnenrei}は組方向に拘わらず同じ結果を得ます。
これらのコマンドは第2引数によって軽微に挙動が変わります。以下の表を参照してください。

\begin{center}
  \begin{booktabular}{cccccc}
    \diagbox[font=\sffamily\gtfamily\small]{第1引数}{第2引数} & \thead{1 -- 9} & \thead{10} & \thead{11 -- 19} & \thead{20 -- } \\
    \hline
    \raisebox{0.7zw}{\texttt{1} {\scriptsize (つ)}} &
    {\ibmfamily \ajnenrei{1}{1}\,\ajnenrei{1}{2}\,\ajnenrei{1}{3}\middots\ajnenrei{1}{8}\,\ajnenrei{1}{9}} &
    {\ibmfamily \ajnenrei{1}{10}} &
    {\ibmfamily \ajnenrei{1}{11}\,\ajnenrei{1}{12}\,\ajnenrei{1}{13}\middots\ajnenrei{1}{18}\,\ajnenrei{1}{19}} &
    {\ibmfamily \ajnenrei{1}{20}\,\ajnenrei{1}{21}\,\ajnenrei{1}{22}\middots\ajnenrei{1}{98}\,\ajnenrei{1}{99}} \\
    \raisebox{0.7zw}{\texttt{2} {\scriptsize (才)}} &
    {\ibmfamily \ajnenrei{2}{1}\,\ajnenrei{2}{2}\,\ajnenrei{2}{3}\middots\ajnenrei{2}{8}\,\ajnenrei{2}{9}} &
    {\ibmfamily \ajnenrei{2}{10}} &
    {\ibmfamily \ajnenrei{2}{11}\,\ajnenrei{2}{12}\,\ajnenrei{2}{13}\middots\ajnenrei{2}{18}\,\ajnenrei{2}{19}} &
    {\ibmfamily \ajnenrei{2}{20}\,\ajnenrei{2}{21}\,\ajnenrei{2}{22}\middots\ajnenrei{2}{98}\,\ajnenrei{2}{99}} \\
    \raisebox{0.7zw}{\texttt{3} {\scriptsize (ツ)}} &
    {\ibmfamily \pbox<t>{\ajnenrei{3}{1}}\,\pbox<t>{\ajnenrei{3}{2}}\,\pbox<t>{\ajnenrei{3}{3}}\middots\pbox<t>{\ajnenrei{3}{8}}\,\pbox<t>{\ajnenrei{3}{9}}} &
    {\ibmfamily \pbox<t>{\ajnenrei{3}{10}}} &
    {\ibmfamily \pbox<t>{\ajnenrei{3}{11}}\,\pbox<t>{\ajnenrei{3}{12}}\,\pbox<t>{\ajnenrei{3}{13}}\middots\pbox<t>{\ajnenrei{3}{18}}\,\pbox<t>{\ajnenrei{3}{19}}} &
    {\ibmfamily \pbox<t>{\ajnenrei{3}{20}}\,\pbox<t>{\ajnenrei{3}{21}}\,\pbox<t>{\ajnenrei{3}{22}}\middots\pbox<t>{\ajnenrei{3}{98}}\,\pbox<t>{\ajnenrei{3}{99}}} \\
    \raisebox{0.7zw}{\texttt{4} {\scriptsize (コ)}} &
    {\ibmfamily \pbox<t>{\ajnenrei{4}{1}}\,\pbox<t>{\ajnenrei{4}{2}}\,\pbox<t>{\ajnenrei{4}{3}}\middots\pbox<t>{\ajnenrei{4}{8}}\,\pbox<t>{\ajnenrei{4}{9}}} &
    {\ibmfamily \pbox<t>{\ajnenrei{4}{10}}} &
    {\ibmfamily \pbox<t>{\ajnenrei{4}{11}}\,\pbox<t>{\ajnenrei{4}{12}}\,\pbox<t>{\ajnenrei{4}{13}}\middots\pbox<t>{\ajnenrei{4}{18}}\,\pbox<t>{\ajnenrei{4}{19}}} &
    {\ibmfamily \pbox<t>{\ajnenrei{4}{20}}\,\pbox<t>{\ajnenrei{4}{21}}\,\pbox<t>{\ajnenrei{4}{22}}\middots\pbox<t>{\ajnenrei{4}{98}}\,\pbox<t>{\ajnenrei{4}{99}}} \\
  \end{booktabular}
\end{center}

これに加えて\cmd{ajnunrei\{4\}}のショートカットとして\cmd{ajKosu}が定義されています。

\newpage
\section{分数：\texorpdfstring{\cmd{ajFrac}}{\textbackslash ajFrac}}

数式中の分数と同じようにすることで文中用の分数を扱うことが出来ます。ただし、この分数はグリフとして登録されているもののみに限ります。

\begin{itemize}
  \item \cmd{ajFrac\{\textlrangle{分子}\}\{\textlrangle{分母}\}}
  \item \cmd{ajFrac*\{\textlrangle{分子}\}\{\textlrangle{分母}\}}
\end{itemize}

例えば以下のように利用されます。

\begin{demosource}
  らんま{\textbackslash}ajFrac\{1\}\{2\}、\% \\
  らんま{\textbackslash}ajFrac*\{1\}\{2\}
  \tcblower
  らんま{\ibmfamily\ajFrac{1}{2}}、%
  らんま{\ibmfamily\ajFrac*{1}{2}}
\end{demosource}

分母は2から12であり、真分数となるような45の分数（各2スタイル）が提供されています。加えて、例外として{\eghostguarded{\ibmfamily\ajFrac{0}{3}}}と{\eghostguarded{\ibmfamily\ajFrac{1}{100}}}があります。
また、約分可能な分数は既約分数のグリフが出力されます。以下の表では灰色網掛けとなっている部分が既約分数となっています。

\noindent
\begin{booktabular}{c|*{10}{c}}
  \diagbox[font=\small\gtfamily]{分母}{分子}
                & \texttt{0}   & \texttt{1}    & \texttt{2}    & \texttt{3}    & \texttt{4}    & \texttt{5}    & \texttt{6}    & \texttt{7}   & \middots[0.2zw] \\
  \hline
  \texttt{2}    & -            & \ajFS{1}{2}   & -             & -             & -             & -             & -             & -            & \multirow{12}{*}{\middots[0.2zw]} \\
  \texttt{3}    & \ajFS{0}{3}  & \ajFS{1}{3}   & \ajFS{2}{3}   & -             & -             & -             & -             & -            &  \\
  \texttt{4}    & -            & \ajFS{1}{4}   & \ajFS*{2}{4}  & \ajFS{3}{4}   & -             & -             & -             & -            &  \\
  \texttt{5}    & -            & \ajFS{1}{5}   & \ajFS{2}{5}   & \ajFS{3}{5}   & \ajFS{4}{5}   & -             & -             & -            &  \\
  \texttt{6}    & -            & \ajFS{1}{6}   & \ajFS*{2}{6}  & \ajFS*{3}{6}  & \ajFS*{4}{6}  & \ajFS{5}{6}   & -             & -            &  \\
  \texttt{7}    & -            & \ajFS{1}{7}   & \ajFS{2}{7}   & \ajFS{3}{7}   & \ajFS{4}{7}   & \ajFS{5}{7}   & \ajFS{6}{7}   & -            &  \\
  \texttt{8}    & -            & \ajFS{1}{8}   & \ajFS*{2}{8}  & \ajFS{3}{8}   & \ajFS*{4}{8}  & \ajFS{5}{8}   & \ajFS*{6}{8}  & \ajFS{7}{8}  &  \\
  \texttt{9}    & -            & \ajFS{1}{9}   & \ajFS{2}{9}   & \ajFS*{3}{9}  & \ajFS{4}{9}   & \ajFS{5}{9}   & \ajFS*{6}{9}  & \ajFS{7}{9}  &  \\
  \texttt{10}   & -            & \ajFS{1}{10}  & \ajFS*{2}{10} & \ajFS{3}{10}  & \ajFS*{4}{10} & \ajFS*{5}{10} & \ajFS*{6}{10} & \ajFS{7}{10} &  \\
  \texttt{11}   & -            & \ajFS{1}{11}  & \ajFS{2}{11}  & \ajFS{3}{11}  & \ajFS{4}{11}  & \ajFS{5}{11}  & \ajFS{6}{11}  & \ajFS{7}{11} &  \\
  \texttt{12}   & -            & \ajFS{1}{12}  & \ajFS*{2}{12} & \ajFS*{3}{12} & \ajFS*{4}{12} & \ajFS{5}{12}  & \ajFS*{6}{12} & \ajFS{7}{12} &  \\
  \texttt{100}  & -            & \ajFS{1}{100} & -             & -             & -             & -             & -             & -            &  \\
\end{booktabular}

\vfill
\hfill
\begin{booktabular}{*{5}{c}|c}
  \middots[0.2zw]     & \texttt{8}    & \texttt{9}     & \texttt{10}    & \texttt{11}   &
    \diagbox[font=\small\gtfamily, dir=NE]{分子}{分母} \\
  \hline
  \multirow{8}{*}{\middots}
      & -             & -             & -              & -             & \texttt{2}   \\[-0.5zw]
      &               &               &                &               & \vdots       \\[-0.5zw]
      & -             & -             & -              & -             & \texttt{8}   \\
      & \ajFS{8}{9}   & -             & -              & -             & \texttt{9}   \\
      & \ajFS*{8}{10} & \ajFS{9}{10}  & -              & -             & \texttt{10}  \\
      & \ajFS{8}{11}  & \ajFS{9}{11}  & \ajFS{10}{11}  & -             & \texttt{11}  \\
      & \ajFS*{8}{12} & \ajFS*{9}{12} & \ajFS*{10}{12} & \ajFS{11}{12} & \texttt{12}  \\
      & -             & -             & -              & -             & \texttt{100} \\
\end{booktabular}

\newpage
\section{矢印：\texorpdfstring{\cmd{ajArrow}}{\textbackslash ajArrow}}

\cmd{ajArrow}を用いることで引数に応じた矢印を出力できます。

加えて、引数に{\optast}や\eghostguarded{\texttt{+}}、\eghostguarded{\texttt{/}}を含まないものは\cmd{ajArrow\textlrangle{\jpincmd{引数}}}として表現することも可能です。

\begin{demosource}
  {\textbackslash}ajArrow\{LeftHand\} 前のページ{\textbackslash}hfill\\
  次のページ {\textbackslash}ajArrowRightHand
  \tcblower
  {\ibmfamily\ajArrow{LeftHand}} 前のページ
  \hfill
  次のページ {\ibmfamily\ajArrowRightHand}
\end{demosource}

\begin{center}
  \begin{booktabular}{cp{17zw}}
    \thead{記号} & \thead{引数\hfill 専用コマンド} \\
    \midrule
    \ajAS{LeftTriangle}         \\
    \ajAS{RightTriangle}        \\
    \ajAS*{LeftTriangle*}       \\
    \ajAS*{RightTriangle*}      \\
    \ajAS{UP}                   \\
    \ajAS{DOWN}                 \\
    \ajAS{LEFT}                 \\
    \ajAS{RIGHT}                \\
    \ajAS{Up}                   \\
    \ajAS{Down}                 \\
    \ajAS{Left}                 \\
    \ajAS{Right}                \\
    \ajAS*{UP*}                 \\
    \ajAS*{DOWN*}               \\
    \ajAS*{LEFT*}               \\
    \ajAS*{RIGHT*}              \\
    \ajAS*{Up+}                 \\
    \ajAS*{Down+}               \\
    \ajAS*{Left+}               \\
    \ajAS*{Right+}              \\
    \ajAS*{UpDown+}             \\
    \ajAS*{LeftRight+}          \\
    \ajAS*{LeftRight*}          \\
    \ajAS*{Up/Down}             \\
    \ajAS*{Down/Up}             \\
    \ajAS*{Left/Right}          \\
    \ajAS*{Right/Left}          \\
    \ajAS*{Down/Up+}            \\
    \ajAS*{Right/Left+}         \\
  \end{booktabular}
  \hfil
  \begin{booktabular}{cp{18zw}}
    \thead{記号} & \thead{引数\hfill 専用コマンド} \\
    \midrule
    \ajAS*{Left/Right*}         \\
    \ajAS*{Right/Left*}         \\
    \ajAS{LeftDouble}           \\
    \ajAS{RightDouble}          \\
    \ajAS{LeftRightDouble}      \\
    \ajAS{LeftUp}               \\
    \ajAS{LeftDown}             \\
    \ajAS{RightUp}              \\
    \ajAS{RightDown}            \\
    \ajAS*{RightUp*}            \\
    \ajAS*{RightDown*}          \\
    \ajAS{UpAngle}              \\
    \ajAS{DownAngle}            \\
    \ajAS{LeftAngle}            \\
    \ajAS{RightAngle}           \\
    \ajAS*{UpAngle*}            \\
    \ajAS*{DownAngle*}          \\
    \ajAS*{LeftAngle*}          \\
    \ajAS*{RightAngle*}         \\
    \ajAS{RightHand}            \\
    \ajAS{UpHand}               \\
    \ajAS{DownHand}             \\
    \ajAS{LeftHand}             \\
    \ajAS{UpScissors}           \\
    \ajAS{DownScissors}         \\
    \ajAS{LeftScissors}         \\
    \ajAS{RightScissors}        \\
  \end{booktabular}
\end{center}

\section{ピクチャー記号：\texorpdfstring{\cmd{ajPICT}}{\textbackslash ajPICT}}

\cmd{ajPICT}を用いることで引数に応じたピクチャー記号を出力できます。
同様に\cmd{\jpincmd{※}}も同じ記号を出力できます。

\begin{demosource}
  カンペキに勝つ{\textbackslash}ajPICT\{Club*\}だろ？\<
  ゴン{\textbackslash}※\{Heart*\}
  \tcblower
  カンペキに勝つ\ajPICT{Club*}だろ？ゴン\※{Heart*}
\end{demosource}

\begin{center}
  \begin{booktabular}{cp{13zw}}
    \thead{記号} & \thead{引数\hfill 専用コマンド} \\
    \midrule
    \ajPS{Club}      \\
    \ajPS*{Club*}    \\
    \ajPS{Heart}     \\
    \ajPS*{Heart*}   \\
    \ajPS{Spade}     \\
    \ajPS*{Spade*}   \\
    \ajPS{Diamond}   \\
    \ajPS*{Diamond*} \\
    \ajPS{→}         \\
    \ajPS{←}         \\
    \ajPS{↑}         \\
    \ajPS{↓}         \\
  \end{booktabular}
  \hfil
  \begin{booktabular}{cp{13zw}}
    \thead{記号} & \thead{引数\hfill 専用コマンド} \\
    \midrule
    \ajPS+{電話}     \\
    \ajPS+{〒}       \\
    \ajPS+{晴}       \\
    \ajPS+{曇}       \\
    \ajPS+{雨}       \\
    \ajPS+{雪}       \\
    \ajPS+{野球}     \\
    \ajPS+{サッカー} \\
    \ajPS+{花}       \\
    \ajPS*+{花*}     \\
    \ajPS+{湯}       \\
  \end{booktabular}
\end{center}

\section{合字(Ligature)：\texorpdfstring{\cmd{ajLig}}{\textbackslash ajLig}}

\cmd{ajLig}を用いることで合字（{\ibmfamily\ajLig{平成}}）や組文字（{\ibmfamily\ajLig{株式会社}}）、囲み文字（{\ibmfamily\ajLig{■例}}）等を得ることが出来ます。

\begin{demosource}
  {\textbackslash}ajLig\{昭和\}、
  {\textbackslash}ajLig\{か゜\}、
  {\textbackslash}ajLig\{km3\}、\\
  {\textbackslash}ajLig\{サンチーム\}、
  {\textbackslash}ajLig\{株式会社\}、
  {\textbackslash}ajLig\{○〒\}。
  \tcblower
  \ajLig{昭和}、
  \ajLig{か゜}、
  \ajLig{km3}、
  \ajLig{サンチーム}、
  \ajLig{株式会社}、
  \ajLig{○〒}。
\end{demosource}

\vfil
\subsection{和文合字、年号、感嘆符・疑問符}

\begin{center}
  \begin{booktabular}{>{\ttfamily\small}p{3zw}c}
    \thead{引数} & \thead{合字} \\
    \midrule
    \ajLigKana{より}   \\
    \ajLigKana{升}     \\
    \ajLigKana{コト}   \\
    \ajLigKana{JIS}    \\
    \ajLigKana*{JAS}    \\
  \end{booktabular}
  \hfil
  \begin{booktabular}{>{\ttfamily\small}C{3zw}cc}
    \thead{引数} & \thead{横組} & \thead{縦組} \\
    \midrule
    \ajLigKumimoji{明治}{0} \\
    \ajLigKumimoji{大正}{0} \\
    \ajLigKumimoji{昭和}{0} \\
    \ajLigKumimoji{平成}{1} \\
    \ajLigKumimoji{令和}{7} \\
  \end{booktabular}
  \hfil
  \begin{booktabular}{>{\ttfamily}C{3zw}c}
    \thead{引数} & \thead{合字} \\
    \midrule
    \ajLigPunc{!!}  \\
    \ajLigPunc{!?}  \\
    \ajLigPunc{?!}  \\
    \ajLigPunc{??}  \\
  \end{booktabular}
  \hfil
  \begin{booktabular}{>{\ttfamily}C{3zw}c}
    \thead{引数} & \thead{合字} \\
    \midrule
    \ajLigPunc{!*}  \\
    \ajLigPunc{!!*} \\
    \ajLigPunc{!?*} \\
  \end{booktabular}
\end{center}

\pkg{otf}では感嘆符・疑問符の後にアキが挿入されますが、\cmd{ajLig}で表現される感嘆符・疑問符の後にアキは挿入されません。

\begin{demosource}
  なんで{\textbackslash}ajLig\{??\}なんで
  \tcblower
  なんで\ajLig{??}なんで
\end{demosource}

\vfil
\subsection{仮名の濁音化・半濁音化、小書き}

濁音には「\,{\ibmfamily \raisebox{-0.2zw}{゛}}\subUCN{309B}」、半濁音には「\,{\ibmfamily \raisebox{-0.2zw}{゜}}\subUCN{309C}」を使用します。

\begin{center}
  \begin{booktabular}{>{\ttfamily\small}lc}
    \thead{引数} & \thead{合字} \\
    \midrule
    \ajLigKana{う゛} \\
    \ajLigKana{ワ゛} \\
    \ajLigKana{ヰ゛} \\
    \ajLigKana{ヱ゛} \\
    \ajLigKana{ヲ゛} \\
  \end{booktabular}\hfil
  \colorbox{\tabularColor}{%
  \begin{booktabular}{>{\ttfamily\small}lc}
    \thead{引数} & \thead{合字} \\
    \midrule
    \ajLigKana{か゜} \\
    \ajLigKana{き゜} \\
    \ajLigKana{く゜} \\
    \ajLigKana{け゜} \\
    \ajLigKana{こ゜} \\
  \end{booktabular}}\hfil
  \begin{booktabular}{>{\ttfamily\small}lc}
    \thead{引数} & \thead{合字} \\
    \midrule
    \ajLigKana{カ゜} \\
    \ajLigKana{キ゜} \\
    \ajLigKana{ク゜} \\
    \ajLigKana{ケ゜} \\
    \ajLigKana{コ゜} \\
  \end{booktabular}\hfil
  \colorbox{\tabularColor}{%
  \begin{booktabular}{>{\ttfamily\small}lc}
    \thead{引数} & \thead{合字} \\
    \midrule
    \ajLigKana{セ゜} \\
    \ajLigKana{ツ゜} \\
    \ajLigKana{ト゜} \\
  \end{booktabular}}\hfil
  \begin{booktabular}{>{\ttfamily\small}lc}
    \thead{引数} & \thead{合字} \\
    \midrule
    \ajLigKana{小か} \\
    \ajLigKana{小け} \\
    \ajLigKana*{小こ} \\
    \ajLigKana*{小コ} \\
    \ajLigKana{小ク} \\
    \ajLigKana{小シ} \\
    \ajLigKana{小ス} \\
  \end{booktabular}\hfil
  \colorbox{\tabularColor}{%
  \begin{booktabular}{>{\ttfamily\small}lc}
    \thead{引数} & \thead{合字} \\
    \midrule
    \ajLigKana{小ト} \\
    \ajLigKana{小ヌ} \\
    \ajLigKana{小ハ} \\
    \ajLigKana{小ヒ} \\
    \ajLigKana{小フ} \\
    \ajLigKana{小ヘ} \\
    \ajLigKana{小ホ} \\
  \end{booktabular}}\hfil
  \begin{booktabular}{>{\ttfamily\small}lc}
    \thead{引数} & \thead{合字} \\
    \midrule
    \ajLigKana{小プ} \\
    \ajLigKana{小ム} \\
    \ajLigKana{小ラ} \\
    \ajLigKana{小リ} \\
    \ajLigKana{小ル} \\
    \ajLigKana{小レ} \\
    \ajLigKana{小ロ} \\
  \end{booktabular}
\end{center}

これらの仮名は濁音・半濁音を組み合わせたり大きさを調整しているわけではなく、そのように構成されたグリフを使用しています。そのため、\cmd{ajLig\{\jpincmd{あ゛}\}}としても「あ」に濁点は与えられません。

\vfil
\subsection{欧文合字}

80の欧文合字があります。
引数となるラテン文字のcaseやスラッシュ(\eghostguarded{\texttt{/}})、ピリオド(\eghostguarded{\texttt{.}})を忘れないようにしてください。

\newcommand{\ajLigLt}[2]{#1 & {\ibmfamily\ajLig{#1}}}

\begin{center}
  \begin{booktabular}{>{\ttfamily\small}lc}
    \thead{引数} & \thead{合字} \\
    \midrule
    \ajLigLt{a.m.}{4}   \\
    \ajLigLt{a/c}{4}    \\
    \ajLigLt{AM}{4}     \\
    \ajLigLt{c.c.}{4}   \\
    \ajLigLt{c/c}{4}    \\
    \ajLigLt{c/o}{4}    \\
    \ajLigLt{cal}{4}    \\
    \ajLigLt{cc}{4}     \\
    \ajLigLt{cm}{4}     \\
    \ajLigLt{cm2}{4}    \\
    \ajLigLt{cm3}{4}    \\
    \ajLigLt{dB}{4}     \\
    \ajLigLt{dl}{4}     \\
    \ajLigLt{dl*}{4}    \\
    \ajLigLt{Dr.}{6}    \\
    \ajLigLt{F}{4}      \\
  \end{booktabular}\hfil
  \colorbox{\tabularColor}{%
  \begin{booktabular}{>{\ttfamily\small}lc}
    \thead{引数} & \thead{合字} \\
    \midrule
    \ajLigLt{FAX}{4}    \\
    \ajLigLt{g}{4}      \\
    \ajLigLt{GB}{4}     \\
    \ajLigLt{H2}{6}     \\
    \ajLigLt{HP}{4}     \\
    \ajLigLt{hPa}{4}    \\
    \ajLigLt{Hz}{4}     \\
    \ajLigLt{in}{4}     \\
    \ajLigLt{Jr.}{6}    \\
    \ajLigLt{K.K.}{4}   \\
    \ajLigLt{KB}{4}     \\
    \ajLigLt{kcal}{4}   \\
    \ajLigLt{kg}{4}     \\
    \ajLigLt{KK.}{4}    \\
    \ajLigLt{KK}{4}     \\
    \ajLigLt{kl}{4}     \\
  \end{booktabular}}\hfil
  \begin{booktabular}{>{\ttfamily\small}lc}
    \thead{引数} & \thead{合字} \\
    \midrule
    \ajLigLt{kl*}{4}    \\
    \ajLigLt{km}{4}     \\
    \ajLigLt{km2}{4}    \\
    \ajLigLt{km3}{4}    \\
    \ajLigLt{l}{4}      \\
    \ajLigLt{l*}{4}     \\
    \ajLigLt{m}{4}      \\
    \ajLigLt{m/m}{4}    \\
    \ajLigLt{m2}{4}     \\
    \ajLigLt{m3}{4}     \\
    \ajLigLt{mb}{4}     \\
    \ajLigLt{MB}{4}     \\
    \ajLigLt{mg}{4}     \\
    \ajLigLt{mho}{6}    \\
    \ajLigLt{microg}{4} \\
    \ajLigLt{microm}{4} \\
  \end{booktabular}\hfil
  \colorbox{\tabularColor}{%
  \begin{booktabular}{>{\ttfamily\small}lc}
    \thead{引数} & \thead{合字} \\
    \midrule
    \ajLigLt{micros}{4} \\
    \ajLigLt{min}{4}    \\
    \ajLigLt{ml}{4}     \\
    \ajLigLt{ml*}{4}    \\
    \ajLigLt{mm}{4}     \\
    \ajLigLt{mm2}{4}    \\
    \ajLigLt{mm3}{4}    \\
    \ajLigLt{ms}{4}     \\
    \ajLigLt{n/m}{4}    \\
    \ajLigLt{No.}{4}    \\
    \ajLigLt{No}{4}     \\
    \ajLigLt{ns}{4}     \\
    \ajLigLt{Nx}{6}     \\
    \ajLigLt{O2}{6}     \\
    \ajLigLt{ohm*}{4}    \\
    \ajLigLt{Ox}{6}     \\
  \end{booktabular}}\hfil
  \begin{booktabular}{>{\ttfamily\small}lc}
    \thead{引数} & \thead{合字} \\
    \midrule
    \ajLigLt{p.m.}{6}   \\
    \ajLigLt{PH}{4}     \\
    \ajLigLt{pH}{6}     \\
    \ajLigLt{PM}{4}     \\
    \ajLigLt{ppb}{6}    \\
    \ajLigLt{ppm}{6}    \\
    \ajLigLt{PR}{4}     \\
    \ajLigLt{ps}{4}     \\
    \ajLigLt{Q2}{6}     \\
    \ajLigLt{sec}{4}    \\
    \ajLigLt{TB}{4}     \\
    \ajLigLt{tel}{4}    \\
    \ajLigLt{Tel}{4}    \\
    \ajLigLt{TEL}{4}    \\
    \ajLigLt{tm}{4}     \\
    \ajLigLt{VS}{4}     \\
  \end{booktabular}
\end{center}

\newpage
\subsection{組文字}

125の仮名組文字と12の漢字組文字があります。

\begin{center}
  \begin{booktabular}{>{\ttfamily\small}p{7zw}*{2}{c}}
    \thead{引数} & \thead{横組} & \thead{縦組} \\
    \midrule
    \ajLigKumimoji{アール}{4}         \\
    \ajLigKumimoji{アト}{5}           \\
    \ajLigKumimoji{アパート}{4}       \\
    \ajLigKumimoji{アルファ}{5}       \\
    \ajLigKumimoji{アンペア}{5}       \\
    \ajLigKumimoji{イニング}{5}       \\
    \ajLigKumimoji{インチ}{4}         \\
    \ajLigKumimoji{ウォン}{5}         \\
    \ajLigKumimoji{ウルシ}{5}         \\
    \ajLigKumimoji{エーカー}{5}       \\
    \ajLigKumimoji{エクサ}{5}         \\
    \ajLigKumimoji{エスクード}{5}     \\
    \ajLigKumimoji{オーム}{5}         \\
    \ajLigKumimoji{オンス}{5}         \\
    \ajLigKumimoji{オントロ}{5}       \\
    \ajLigKumimoji{カイリ}{5}         \\
    \ajLigKumimoji{カップ}{5}         \\
    \ajLigKumimoji{カラット}{5}       \\
    \ajLigKumimoji{ガル}{6}           \\
    \ajLigKumimoji{カロリー}{4}       \\
    \ajLigKumimoji{ガロン}{5}         \\
    \ajLigKumimoji{ガンマ}{5}         \\
    \ajLigKumimoji{ギガ}{5}           \\
    \ajLigKumimoji{ギニー}{5}         \\
    \ajLigKumimoji{キュリー}{5}       \\
    \ajLigKumimoji{ギルダー}{5}       \\
    \ajLigKumimoji{キロ}{4}           \\
    \ajLigKumimoji{キログラム}{4}     \\
    \ajLigKumimoji{キロメートル}{4}   \\
    \ajLigKumimoji{キロリットル}{5}   \\
    \ajLigKumimoji{キロワット}{5}     \\
    \ajLigKumimoji{グスーム}{5}       \\
    \ajLigKumimoji{グラム}{4}         \\
    \ajLigKumimoji{グラムトン}{5}     \\
  \end{booktabular}\hfil
  \colorbox{\tabularColor}{%
  \begin{booktabular}{>{\ttfamily\small}p{7zw}*{2}{c}}
    \thead{引数} & \thead{横組} & \thead{縦組} \\
    \midrule
    \ajLigKumimoji{クルサード}{5}     \\
    \ajLigKumimoji{クルゼイロ}{5}     \\
    \ajLigKumimoji{グレイ}{6}         \\
    \ajLigKumimoji{クローネ}{5}       \\
    \ajLigKumimoji{ケース}{5}         \\
    \ajLigKumimoji{コーポ}{4}         \\
    \ajLigKumimoji{コルナ}{5}         \\
    \ajLigKumimoji{サイクル}{5}       \\
    \ajLigKumimoji{さじ}{5}           \\
    \ajLigKumimoji{サンチーム}{5}     \\
    \ajLigKumimoji{シーベルト}{6}     \\
    \ajLigKumimoji{シェケル}{6}       \\
    \ajLigKumimoji{ジュール}{6}       \\
    \ajLigKumimoji{シリング}{5}       \\
    \ajLigKumimoji{センチ}{4}         \\
    \ajLigKumimoji{セント}{4}         \\
    \ajLigKumimoji{ダース}{5}         \\
    \ajLigKumimoji{デカ}{5}           \\
    \ajLigKumimoji{デシ}{5}           \\
    \ajLigKumimoji{デシベル}{6}       \\
    \ajLigKumimoji{テラ}{5}           \\
    \ajLigKumimoji{ドット}{6}         \\
    \ajLigKumimoji{ドラクマ}{5}       \\
    \ajLigKumimoji{ドル}{4}           \\
    \ajLigKumimoji{トン}{4}           \\
    \ajLigKumimoji{ナノ}{5}           \\
    \ajLigKumimoji{ノット}{5}         \\
    \ajLigKumimoji{パーセント}{4}     \\
    \ajLigKumimoji{バーツ}{5}         \\
    \ajLigKumimoji{バーレル}{5}       \\
    \ajLigKumimoji{ハイツ}{4}         \\
    \ajLigKumimoji{バイト}{6}         \\
    \ajLigKumimoji{パスカル}{5}       \\
    \ajLigKumimoji{バレル}{5}         \\
  \end{booktabular}}\hfil
  \begin{booktabular}{>{\ttfamily\small}p{7zw}*{2}{c}}
    \thead{引数} & \thead{横組} & \thead{縦組} \\
    \midrule
    \ajLigKumimoji{ピアストル}{5}     \\
    \ajLigKumimoji{ピクル}{5}         \\
    \ajLigKumimoji{ピコ}{5}           \\
    \ajLigKumimoji{ビット}{6}         \\
    \ajLigKumimoji{ビル}{4}           \\
    \ajLigKumimoji{ファラッド}{5}     \\
    \ajLigKumimoji{ファラド}{5}       \\
    \ajLigKumimoji{フィート}{5}       \\
    \ajLigKumimoji{フェムト}{5}       \\
    \ajLigKumimoji{ブッシェル}{5}     \\
    \ajLigKumimoji{フラン}{5}         \\
    \ajLigKumimoji{ページ}{4}         \\
    \ajLigKumimoji{ベータ}{5}         \\
    \ajLigKumimoji{ヘクタール}{4}     \\
    \ajLigKumimoji{ヘクト}{5}         \\
    \ajLigKumimoji{ヘクトパスカル}{5} \\
    \ajLigKumimoji{ベクレル}{6}       \\
    \ajLigKumimoji{ペセタ}{5}         \\
    \ajLigKumimoji{ペソ}{5}           \\
    \ajLigKumimoji{ペタ}{5}           \\
    \ajLigKumimoji{ペニヒ}{5}         \\
    \ajLigKumimoji{ヘルツ}{4}         \\
    \ajLigKumimoji{ペンス}{5}         \\
    \ajLigKumimoji{ボー}{6}           \\
    \ajLigKumimoji{ホール}{5}         \\
    \ajLigKumimoji{ホーン}{4}         \\
    \ajLigKumimoji{ポイント}{5}       \\
    \ajLigKumimoji{ボルト}{5}         \\
    \ajLigKumimoji{ホン}{5}           \\
    \ajLigKumimoji{ポンド}{5}         \\
    \ajLigKumimoji{マイクロ}{5}       \\
    \ajLigKumimoji{マイル}{5}         \\
    \ajLigKumimoji{マッハ}{5}         \\
    \ajLigKumimoji{マルク}{5}         \\
  \end{booktabular}
\end{center}

\begin{center}
  \colorbox{\tabularColor}{%
  \begin{booktabular}{>{\ttfamily\small}p{7zw}*{2}{c}}
    \thead{引数} & \thead{横組} & \thead{縦組} \\
    \midrule
    \ajLigKumimoji{マンション}{4}     \\
    \ajLigKumimoji{ミクロン}{5}       \\
    \ajLigKumimoji{ミリ}{4}           \\
    \ajLigKumimoji{ミリバール}{4}     \\
    \ajLigKumimoji{メートル}{4}       \\
    \ajLigKumimoji{メガ}{5}           \\
    \ajLigKumimoji{メガトン}{5}       \\
    \ajLigKumimoji{ヤード}{4}         \\
    \ajLigKumimoji{ヤール}{5}         \\
    \ajLigKumimoji{ユーロ}{5}         \\
    \ajLigKumimoji{ユアン}{5}         \\
    \ajLigKumimoji{ラド}{5}           \\
  \end{booktabular}}\hfil
  \begin{booktabular}{>{\ttfamily\small}p{7zw}*{2}{c}}
    \thead{引数} & \thead{横組} & \thead{縦組} \\
    \midrule
    \ajLigKumimoji{ランド}{6}         \\
    \ajLigKumimoji{リットル}{4}       \\
    \ajLigKumimoji{リラ}{5}           \\
    \ajLigKumimoji{リンギット}{6}     \\
    \ajLigKumimoji{ルーブル}{5}       \\
    \ajLigKumimoji{ルクス}{5}         \\
    \ajLigKumimoji{ルピー}{5}         \\
    \ajLigKumimoji{ルピア}{5}         \\
    \ajLigKumimoji{レム}{5}           \\
    \ajLigKumimoji{レントゲン}{5}     \\
    \ajLigKumimoji{ワット}{4}         \\
  \end{booktabular}\hfil
  \colorbox{\tabularColor}{%
  \begin{booktabular}{>{\ttfamily\small}p{7zw}*{2}{c}}
    \thead{引数} & \thead{横組} & \thead{縦組} \\
    \midrule
    \ajLigKumimoji{株式会社}{4}       \\
    \ajLigKumimoji{有限会社}{4}       \\
    \ajLigKumimoji{財団法人}{4}       \\
    \ajLigKumimoji{医療法人}{5}       \\
    \ajLigKumimoji{学校法人}{5}       \\
    \ajLigKumimoji{協同組合}{5}       \\
    \ajLigKumimoji{共同組合}{5}       \\
    \ajLigKumimoji{合資会社}{5}       \\
    \ajLigKumimoji{合名会社}{5}       \\
    \ajLigKumimoji{社団法人}{5}       \\
    \ajLigKumimoji{宗教法人}{5}       \\
    \ajLigKumimoji{郵便番号}{5}       \\
  \end{booktabular}}
\end{center}

% \newpage
\subsubsection{変種のあるもの}

組文字には特に3字のものに対して変種となるグリフがあります。これらを利用するには、引数の末に{\optast}を与えます。この変種グリフは1文字の部分が中央揃えになります。

\begin{demosource}
  {\textbackslash}ajLig\{\jpincmd{センチ}\}{\textbackslash},\\
  {\textbackslash}ajLig\{\jpincmd{センチ}*\}
  \tcblower
  \ibmfamily\LARGE
  \ajLig{センチ}\,
  \ajLig{センチ*}
\end{demosource}

\begin{center}
  \begin{booktabular}{>{\ttfamily\small}ccccc}
    \thead{引数} & \thead{横組} & \thead{横組\optast} & \thead{縦組} & \thead{縦組\optast} \\
    \midrule
    \ajLigKumimoji*{センチ}{4} \\
    \ajLigKumimoji*{グラム}{4} \\
    \ajLigKumimoji*{アール}{4} \\
    \ajLigKumimoji*{ワット}{4} \\
    \ajLigKumimoji*{セント}{4} \\
    \ajLigKumimoji*{ページ}{4} \\
  \end{booktabular}
  \hfil
  \begin{booktabular}{*{2}{>{\ttfamily\small}p{4zw}cccc}}
    \thead{引数} & \thead{横組} & \thead{横組\optast} & \thead{縦組} & \thead{縦組\optast} \\
    \midrule
    \ajLigKumimoji*{インチ}{4} \\
    \ajLigKumimoji*{ヤード}{4} \\
    \ajLigKumimoji*{ヘルツ}{4} \\
    \ajLigKumimoji*{ホーン}{4} \\
    \ajLigKumimoji*{コーポ}{4} \\
    \ajLigKumimoji*{ハイツ}{4} \\
  \end{booktabular}
\end{center}

3文字の組文字はこれら以外にも多くありますが、この変種に対応するグリフは{\AdobeJapan[4]}以前に収録されたグリフであり、それ以降の追補では集録されていないため存在しません。

加えて、\cmd{ajLig\{\jpincmd{オングストローム}\}}とその変種があります。この組文字は1文字幅で表現するには非常に文字が細かいため、これを2つに分けた\cmd{ajLig\{\jpincmd{オントロ}\}}（{\ibmfamily\ajLig{オントロ}}）と\cmd{ajLig\{\jpincmd{グスーム}\}}（{\ibmfamily\ajLig{グスーム}}）があります。オングストロームの変種では、これらを連結した2文字幅を出力します。

\begin{center}
  %% オングストロームは AJ1-4
  \begin{booktabular}{>{\ttfamily\small}l*{4}{>{\ibmfamily}{c}}}
    \thead{引数} & \thead{横組} & \thead{横組\optast} & \thead{縦組} & \thead{縦組\optast} \\
    \midrule
    オングストローム &
    \ajLig{オングストローム} &
    \ajLig{オングストローム*} &
    \raisebox{0.1zw}{\parbox<t>{1zw}{\ajLig{オングストローム}}} &
    \parbox<t>{2zw}{\ajLig{オングストローム*}} \\
  \end{booktabular}
\end{center}

\newpage
\subsection{囲み文字}

143の囲み文字があります。
引数に利用される記号はそれぞれ「○\subUCN{25CB}」、「●\subUCN{25CF}」、「□\subUCN{25A1}」、「■\subUCN{25A0}」、「◇\subUCN{25C7}」、「◆\subUCN{25C6}」、丸括弧は半角丸括弧です。

\newcommand{\ajLigKakomi}[1]{#1 & {\ibmfamily \ajLig{#1}}}

\begin{center}
  \begin{booktabular}{>{\ttfamily\small}p{3zw}c}
    \thead{引数} & \thead{合字} \\
    \midrule
    \ajLigKakomi{○問} \\
    \ajLigKakomi{○答} \\
    \ajLigKakomi{○例} \\
    \ajLigKakomi{○〒} \\
    \ajLigKakomi{○上} \\
    \ajLigKakomi{○下} \\
    \ajLigKakomi{○中} \\
    \ajLigKakomi{○企} \\
    \ajLigKakomi{○休} \\
    \ajLigKakomi{○優} \\
    \ajLigKakomi{○写} \\
    \ajLigKakomi{○出} \\
    \ajLigKakomi{○副} \\
    \ajLigKakomi{○労} \\
    \ajLigKakomi{○医} \\
    \ajLigKakomi{○協} \\
    \ajLigKakomi{○印} \\
    \ajLigKakomi{○右} \\
    \ajLigKakomi{○名} \\
    \ajLigKakomi{○基} \\
    \ajLigKakomi{○増} \\
    \ajLigKakomi{○夜} \\
    \ajLigKakomi{○大} \\
    \ajLigKakomi{○女} \\
  \end{booktabular}\hfil
  \colorbox{\tabularColor}{%
  \begin{booktabular}{>{\ttfamily\small}p{3zw}c}
    \thead{引数} & \thead{合字} \\
    \midrule
    \ajLigKakomi{○学} \\
    \ajLigKakomi{○宗} \\
    \ajLigKakomi{○小} \\
    \ajLigKakomi{○左} \\
    \ajLigKakomi{○年} \\
    \ajLigKakomi{○控} \\
    \ajLigKakomi{○有} \\
    \ajLigKakomi{○株} \\
    \ajLigKakomi{○標} \\
    \ajLigKakomi{○欠} \\
    \ajLigKakomi{○正} \\
    \ajLigKakomi{○注} \\
    \ajLigKakomi{○済} \\
    \ajLigKakomi{○減} \\
    \ajLigKakomi{○特} \\
    \ajLigKakomi{○男} \\
    \ajLigKakomi{○監} \\
    \ajLigKakomi{○社} \\
    \ajLigKakomi{○祝} \\
    \ajLigKakomi{○禁} \\
    \ajLigKakomi{○秘} \\
    \ajLigKakomi{○調} \\
    \ajLigKakomi{○財} \\
    \ajLigKakomi{○資} \\
  \end{booktabular}}\hfil
  \begin{booktabular}{>{\ttfamily\small}p{3zw}c}
    \thead{引数} & \thead{合字} \\
    \midrule
    \ajLigKakomi{○適} \\
    \ajLigKakomi{○電} \\
    \ajLigKakomi{○項} \\
    \ajLigKakomi{●問} \\
    \ajLigKakomi{●答} \\
    \ajLigKakomi{●例} \\
    \ajLigKakomi{□問}  \\
    \ajLigKakomi{□答}  \\
    \ajLigKakomi{□例}  \\
    \ajLigKakomi{□勝}  \\
    \ajLigKakomi{□印}  \\
    \ajLigKakomi{□負}  \\
    \ajLigKakomi{■問}  \\
    \ajLigKakomi{■答}  \\
    \ajLigKakomi{■例}  \\
    \ajLigKakomi{◇問}    \\
    \ajLigKakomi{◇答}    \\
    \ajLigKakomi{◇例}    \\
    \ajLigKakomi{◇HV}    \\
    \ajLigKakomi{◇MV}    \\
    \ajLigKakomi{◇News}  \\
    \ajLigKakomi{◇PV}    \\
    \ajLigKakomi{◇S1}    \\
    \ajLigKakomi{◇S2}    \\
  \end{booktabular}\hfil
  \colorbox{\tabularColor}{%
  \begin{booktabular}{>{\ttfamily\small}p{3zw}c}
    \thead{引数} & \thead{合字} \\
    \midrule
    \ajLigKakomi{◇S3}    \\
    \ajLigKakomi{◇SS}    \\
    \ajLigKakomi{◇デ}    \\
    \ajLigKakomi{◇ほか}  \\
    \ajLigKakomi{◇二}    \\
    \ajLigKakomi{◇交}    \\
    \ajLigKakomi{◇再}    \\
    \ajLigKakomi{◇前}    \\
    \ajLigKakomi{◇劇}    \\
    \ajLigKakomi{◇双}    \\
    \ajLigKakomi{◇司}    \\
    \ajLigKakomi{◇声}    \\
    \ajLigKakomi{◇多}    \\
    \ajLigKakomi{◇天}    \\
    \ajLigKakomi{◇後}    \\
    \ajLigKakomi{◇手}    \\
    \ajLigKakomi{◇文}    \\
    \ajLigKakomi{◇新}    \\
    \ajLigKakomi{◇映}    \\
    \ajLigKakomi{◇株}    \\
    \ajLigKakomi{◇気}    \\
    \ajLigKakomi{◇立}    \\
    \ajLigKakomi{◇終}    \\
    \ajLigKakomi{◇解}    \\
  \end{booktabular}}\hfil
  \begin{booktabular}{>{\ttfamily\small}p{3zw}c}
    \thead{引数} & \thead{合字} \\
    \midrule
    \ajLigKakomi{◆問}    \\
    \ajLigKakomi{◆答}    \\
    \ajLigKakomi{◆例}    \\
    \ajLigKakomi{(問)}  \\
    \ajLigKakomi{(答)}  \\
    \ajLigKakomi{(例)}  \\
    \ajLigKakomi{(代)}  \\
    \ajLigKakomi{(企)}  \\
    \ajLigKakomi{(休)}  \\
    \ajLigKakomi{(労)}  \\
    \ajLigKakomi{(協)}  \\
    \ajLigKakomi{(合)}  \\
    \ajLigKakomi{(名)}  \\
    \ajLigKakomi{(呼)}  \\
    \ajLigKakomi{(営)}  \\
    \ajLigKakomi{(学)}  \\
    \ajLigKakomi{(有)}  \\
    \ajLigKakomi{(株)}  \\
    \ajLigKakomi{(注)}  \\
    \ajLigKakomi{(特)}  \\
    \ajLigKakomi{(監)}  \\
    \ajLigKakomi{(社)}  \\
    \ajLigKakomi{(祝)}  \\
    \ajLigKakomi{(祭)}  \\
  \end{booktabular}\hfil
  \colorbox{\tabularColor}{%
  \begin{booktabular}{>{\ttfamily\small}p{3zw}c}
    \thead{引数} & \thead{合字} \\
    \midrule
    \ajLigKakomi{(自)}  \\
    \ajLigKakomi{(至)}  \\
    \ajLigKakomi{(財)}  \\
    \ajLigKakomi{(資)}  \\
    \ajLigKakomi{□:A}   \\
    \ajLigKakomi{□:AM}  \\
    \ajLigKakomi{□:AS}  \\
    \ajLigKakomi{□:B}   \\
    \ajLigKakomi{□:BEL} \\
    \ajLigKakomi{□:C}   \\
    \ajLigKakomi{□:CL}  \\
    \ajLigKakomi{□:D}   \\
    \ajLigKakomi{□:E}   \\
    \ajLigKakomi{□:F}   \\
    \ajLigKakomi{□:KCL} \\
    \ajLigKakomi{□:ゴ}  \\
    \ajLigKakomi{□:ミ}  \\
    \ajLigKakomi{□:段}  \\
    \ajLigKakomi{□:終}  \\
    \ajLigKakomi{■◇}  \\
    \ajLigKakomi{△!}  \\
    \ajLigKakomi{▽▽}  \\
    \ajLigKakomi{▽〒} \\
  \end{booktabular}}
\end{center}

\cmd{ajMaruKakuAlph\{14\}}（{\ibmfamily\ajMaruKakuAlph{14}}）と\cmd{ajLig\{◇News\}}（{\ibmfamily\ajLig{◇News}}）は見た目が変わらないグリフですが、これらは異なるグリフです。
加えて、\texttt{CID+20504}（{\ibmfamily\CID{20504}}）は提供されていません。ただし、これは\cmd{ajMaruAlph\{8\}}（{\ibmfamily\ajMaruAlph{8}}）とは異なるグリフです。

一部、U-Pressから導入されたグリフがあります。これらは{\Unicode}における私用領域に当てられていますが、PDF上では私用領域にマッピングされていないことに注意してください。
次の46のグリフが対象です：｛{\ibmfamily\□{:A}, \□{:B}, \□{:C}, \□{:D}, \□{:E}, \□{:F}, \□{:終}, \□{:CL}, \□{:KCL}, \□{:BEL}, \□{:AS}, \□{:AM}, \□{:段}, \□{:ゴ}, \□{:ミ}, \○{年}, \◇{News}, \◇{天}, \◇{再}, \◇{新}, \◇{映}, \◇{声}, \◇{前}, \◇{後}, \◇{終}, \◇{立}, \◇{交}, \◇{ほか}, \◇{劇}, \◇{司}, \◇{解}, \◇{株}, \◇{気}, \◇{二}, \◇{多}, \◇{文}, \◇{手}, \◇{PV}, \◇{MV}, \◇{双}, \◇{SS}, \◇{S1}, \◇{S2}, \◇{S3}, \◇{デ}, \◇{HV}.}｝

\newpage
\section{コントロールシンボル}

\pkg{ajmacros}では、囲み文字や合字を``コントロールシンボル''と呼ばれる方法で直接表現することが出来ます。
基本形は\cmd{\textlrangle{記号}\textlrangle{文字}}とします。
この\eghostguarded{\textlrangle{記号}}は囲み文字のための9種と仮名の濁音化・半濁音化、小書きのための3種があります。

ここで利用される記号には類似した文字が多数あるため、以下の表を参考に記号に利用する文字に注意してください。

\begin{center}
  \begin{booktabular}{>{\ttfamily}C{5zw}C{10zw}}
    \thead{\textlrangle{記号}} & \thead{{\Unicode}} \\
    \midrule
    ○            & \UCN{25CB}      \\
    ●            & \UCN{25CF}      \\
    □            & \UCN{25A1}      \\
    ■            & \UCN{25A0}      \\
    ◇            & \UCN{25C7}      \\
    ◆            & \UCN{25C6}      \\
  \end{booktabular}
  \hfil
  \begin{booktabular}{>{\ttfamily}C{5zw}C{10zw}}
    \thead{\textlrangle{記号}} & \thead{{\Unicode}} \\
    \midrule
    （{\rmfamily,} ）       & \UCN{FF08}, \UCN{FF09} \\
    △            & \UCN{25B3}      \\
    ▽            & \UCN{25BD}      \\
    ゛           & \UCN{309B}      \\
    ゜           & \UCN{309C}      \\
    ！           & \UCN{FF01}      \\
  \end{booktabular}
\end{center}

実際には以下のように利用します。

\begin{demosource}
  {\textbackslash}○a、{\textbackslash}●B、{\textbackslash}□う、{\textbackslash}■ヱ、{\textbackslash}◇木、{\textbackslash}◆か、{\textbackslash}（問）\\
  {\textbackslash}△!、{\textbackslash}゛\{\jpincmd{ワ}\}、{\textbackslash}゜\{\jpincmd{セ}\}、{\textbackslash}！\hspace{-0.4zw}\{\jpincmd{ヌ}\}
  \tcblower
  \ibmfamily
  \○a、\●B、\□う、\■ヱ、\◇木、\◆か、\（問）\\
  \△!、\゛{ワ}、\゜{セ}、\！{ヌ}
\end{demosource}

本ドキュメントは{\upLaTeX}で作成しているため仮名のコントロールシンボルでは波括弧で囲っていますが、{\pLaTeX}では囲わなくても同じ結果を得ることが出来ます。（\cmd{kcatcode}を変更することでも対応できます：\href{https://texwiki.texjp.org/?OTF#:~:text=uplatex%20%E3%81%A7%E5%87%A6%E7%90%86%E3%81%99%E3%82%8B%E5%A0%B4%E5%90%88%E3%81%AB%20%5Ckcatcode%60%E3%82%9C%3D18%20%E3%82%92%E8%BF%BD%E5%8A%A0%E3%81%97%E3%81%A6%E3%81%84%E3%81%BE%E3%81%99%E3%80%82%20%E3%81%93%E3%82%8C%E3%82%92%E8%BF%BD%E5%8A%A0%E3%81%97%E3%81%AA%E3%81%84%E5%A0%B4%E5%90%88%E3%81%AF%20uplatex%20%E3%81%A7%E3%82%A8%E3%83%A9%E3%83%BC%E3%81%8C%E7%99%BA%E7%94%9F%E3%81%97%E3%81%BE%E3%81%99%E3%80%82%20upTeX%20%E3%81%A7%E3%81%AE%20kcatcode%20%E3%81%AE%E6%97%A2%E5%AE%9A%E5%80%A4%E3%81%8C%20pTeX%20%E3%81%8B%E3%82%89%E5%A4%89%E6%9B%B4%E3%81%95%E3%82%8C%E3%81%A6%E3%81%84%E3%82%8B%E3%82%88%E3%81%86%E3%81%A7%E3%81%99%E3%80%82}{{\TeXWiki} -- OTF}）

特に、複数文字の引数を採るコントロールシンボルの場合は波括弧で囲う必要があります。以下で示す\cmd{◇ほか}の例では波括弧で囲っていないため\cmd{◇ほ}と解釈されてしまっています。

\begin{demosource}
  {\textbackslash}◇ほか、{\textbackslash}◇\{\jpincmd{ほか}\}。\\
  {\textbackslash}□:ゴ、{\textbackslash}□\{:ゴ\}。
  \% \warichu{\rmfamily\cmd{□:}{\inhibitkanjiskip}\,に相当するグリフは存在しないがエラーにならない}
  \\
  {\textbackslash}◇News、{\textbackslash}◇\{News\}。
  \tcblower
  \ibmfamily
  \◇ほか、\◇{ほか}。\\
  \□:ゴ、\□{:ゴ}。\\
  \◇News、\◇{News}。
\end{demosource}

仮名を濁音化・半濁音化する場合、\cmd{ajLig}で提供される引数と同じように、任意のすべての仮名の濁音化・半濁音化、小書きを表現できるわけではありません。以下の表の通り、取りうる引数は{\Unicode}にあるが{\ShiftJIS}に含まれない文字に限られます。

\begin{center}
  \begin{booktabular}{cl}
    \thead{濁点}   & うワヰヱヲ \\
    \thead{半濁点} & かきくけこカキクケコセツト \\
    \thead{小書き} & かけこクコシストハヒフヘホヌプムラリルレロ \\
  \end{booktabular}
\end{center}

そのため、例えば「あ」に濁点を与えた文字を期待して\cmd{゛\!\{\jpincmd{あ}\}}としても何も出力されません。（濁点・半濁点の結合文字を利用したとしても、NFC正規化が行われるため存在しないグリフを参照することとなり、やはり何も出力されません）

\newpage
\section{\texorpdfstring{\pkg{otf}}{otfパッケージ}から提供されるコマンド}

\pkg{ajmacros}ではなく\pkg{otf}から提供されている記号用コマンドがあります。
これらには、すでに示した\cmd{ajArrow}、\cmd{ajPICT}、\cmd{ajLig}から提供されているコマンドと同様のコマンドが別名で定義されているものも多くあります。

ここでは、すでに示されたコマンド名を並記する形でも\pkg{otf}から提供されるコマンドを紹介します。

\subsection{矢印}

\begin{center}
  \begin{booktabular}{>{\ibmfamily}lp{10zw}*{2}{p{13zw}}>{\ibmfamily}l}
    \ajUpBArrow        & \cmd{ajUpBArrow}       & \cmd{ajArrow\{UP*\}}            & --                         & \ajUpBArrow        \\
    \ajDownBArrow      & \cmd{ajDownBArrow}     & \cmd{ajArrow\{DOWN*\}}          & --                         & \ajDownBArrow      \\
    \ajRightBArrow     & \cmd{ajRightBArrow}    & \cmd{ajArrow\{RIGHT*\}}         & --                         & \ajRightBArrow     \\
    \ajLeftBArrow      & \cmd{ajLeftBArrow}     & \cmd{ajArrow\{LEFT*\}}          & --                         & \ajLeftBArrow      \\
    \ajUpWArrow        & \cmd{ajUpWArrow}       & \cmd{ajArrow\{Up\}}             & \cmd{ajArrowUp}            & \ajUpWArrow        \\
    \ajDownWArrow      & \cmd{ajDownWArrow}     & \cmd{ajArrow\{Down\}}           & \cmd{ajArrowDown}          & \ajDownWArrow      \\
    \ajRightWArrow     & \cmd{ajRightWArrow}    & \cmd{ajArrow\{Right\}}          & \cmd{ajArrowRight}         & \ajRightWArrow     \\
    \ajLeftWArrow      & \cmd{ajLeftWArrow}     & \cmd{ajArrow\{Left\}}           & \cmd{ajArrowLeft}          & \ajLeftWArrow      \\
    \ajLeftUpArrow     & \cmd{ajLeftUpArrow}    & \cmd{ajArrow\{LeftUp\}}         & \cmd{ajArrowLeftUp}        & \ajLeftUpArrow     \\
    \ajRightUpArrow    & \cmd{ajRightUpArrow}   & \cmd{ajArrow\{RightUp\}}        & \cmd{ajArrowRightUp}       & \ajRightUpArrow    \\
    \ajRightDownArrow  & \cmd{ajRightDownArrow} & \cmd{ajArrow\{RightDown\}}      & \cmd{ajArrowRightDown}     & \ajRightDownArrow  \\
    \ajLeftDownArrow   & \cmd{ajLeftDownArrow}  & \cmd{ajArrow\{LeftDown\}}       & \cmd{ajArrowLeftDown}      & \ajLeftDownArrow   \\
    \ajUpScissors      & \cmd{ajUpScissors}     & \cmd{ajArrow\{UpScissors\}}     & \cmd{ajArrowUpScissors}    & \ajUpScissors      \\
    \ajDownScissors    & \cmd{ajDownScissors}   & \cmd{ajArrow\{DownScissors\}}   & \cmd{ajArrowDownScissors}  & \ajDownScissors    \\
    \ajRightScissors   & \cmd{ajRightScissors}  & \cmd{ajArrow\{RightScissors\}}  & \cmd{ajArrowRightScissors} & \ajRightScissors   \\
    \ajLeftScissors    & \cmd{ajLeftScissors}   & \cmd{ajArrow\{LeftScissors\}}   & \cmd{ajArrowLeftScissors}  & \ajLeftScissors    \\
  \end{booktabular}
\end{center}

手で方向を指定するグリフは\cmd{ajArrow}と\cmd{aPICT}の両方があります。

\begin{center}
  \begin{booktabular}{>{\ibmfamily}lp{10zw}*{2}{p{13zw}}>{\ibmfamily}l}
    \ajUpHand          & \cmd{ajUpHand}         & \cmd{ajArrow\{UpHand\}}         & \cmd{ajArrowUpHand}        & \ajUpHand          \\
                       &                        & \cmd{ajPICT\{↑\}}               & \cmd{ajPICT↑}              &                    \\
    \ajDownHand        & \cmd{ajDownHand}       & \cmd{ajArrow\{DownHand\}}       & \cmd{ajArrowDownHand}      & \ajDownHand        \\
                       &                        & \cmd{ajPICT\{↓\}}               & \cmd{ajPICT↓}              &                    \\
    \ajRightHand       & \cmd{ajRightHand}      & \cmd{ajArrow\{RightHand\}}      & \cmd{ajArrowRightHand}     & \ajRightHand       \\
                       &                        & \cmd{ajPICT\{→\}}               & \cmd{ajPICT→}              &                    \\
    \ajLeftHand        & \cmd{ajLeftHand}       & \cmd{ajArrow\{LeftHand\}}       & \cmd{ajArrowLeftHand}      & \ajLeftHand        \\
                       &                        & \cmd{ajPICT\{←\}}               & \cmd{ajPICT←}              &                    \\
  \end{booktabular}
\end{center}

\subsection{ピクチャー記号}

\begin{center}
  \begin{booktabular}{>{\ibmfamily}l*{3}{p{12zw}}>{\ibmfamily}l}
    \ajClub          & \cmd{ajClub}           & \cmd{ajPICT\{Club*\}}          & --                          & \ajClub          \\
    \ajvarClub       & \cmd{ajvarClub}        & \cmd{ajPICT\{Club\}}           & \cmd{ajPICTClub}            & \ajvarClub       \\
    \ajSpade         & \cmd{ajSpade}          & \cmd{ajPICT\{Spade*\}}         & --                          & \ajSpade         \\
    \ajvarSpade      & \cmd{ajvarSpade}       & \cmd{ajPICT\{Spade\}}          & \cmd{ajPICTSpade}           & \ajvarSpade      \\
    \ajvarHeart      & \cmd{ajvarHeart}       & \cmd{ajPICT\{Heart*\}}         & --                          & \ajvarHeart      \\
    \ajHeart         & \cmd{ajHeart}          & \cmd{ajPICT\{Heart\}}          & \cmd{ajPICTHeart}           & \ajHeart         \\
    \ajvarDiamond    & \cmd{ajvarDiamond}     & \cmd{ajPICT\{Diamond*\}}       & --                          & \ajvarDiamond    \\
    \ajDiamond       & \cmd{ajDiamond}        & \cmd{ajPICT\{Diamond\}}        & \cmd{ajPICTDiamond}         & \ajDiamond       \\
    \ajBlackFlorette & \cmd{ajBlackFlorette}  & \cmd{ajPICT\{\jpincmd{花}*\}}  & --                          & \ajBlackFlorette \\
    \ajWhiteFlorette & \cmd{ajWhiteFlorette}  & \cmd{ajPICT\{\jpincmd{花}\}}   & \cmd{ajPICT\jpincmd*{花}}   & \ajWhiteFlorette \\
    \ajPostal        & \cmd{ajPostal}         & \cmd{ajPICT\{〒\}}             & \cmd{ajPICT{\inhibitkanjiskip}〒} & \ajPostal        \\
    \ajvarPostal     & \cmd{ajvarPostal}      & --                             & --                          & \ajvarPostal     \\
    \ajPhone         & \cmd{ajPhone}          & \cmd{ajPICT\{\jpincmd{電話}\}} & \cmd{ajPICT\jpincmd*{電話}} & \ajPhone         \\
    \ajSun           & \cmd{ajSun}            & \cmd{ajPICT\{\jpincmd{晴}\}}   & \cmd{ajPICT\jpincmd*{晴}}   & \ajSun           \\
    \ajCloud         & \cmd{ajCloud}          & \cmd{ajPICT\{\jpincmd{曇}\}}   & \cmd{ajPICT\jpincmd*{曇}}   & \ajCloud         \\
    \ajUmbrella      & \cmd{ajUmbrella}       & \cmd{ajPICT\{\jpincmd{雨}\}}   & \cmd{ajPICT\jpincmd*{雨}}   & \ajUmbrella      \\
    \ajSnowman       & \cmd{ajSnowman}        & \cmd{ajPICT\{\jpincmd{雪}\}}   & \cmd{ajPICT\jpincmd*{雪}}   & \ajSnowman       \\
    \ajBall          & \cmd{ajBall}           & \cmd{ajPICT\{\jpincmd{野球}\}} & \cmd{ajPICT\jpincmd*{野球}} & \ajBall          \\
    \ajHotSpring     & \cmd{ajHotSpring}      & \cmd{ajPICT\{\jpincmd{温泉}\}} & \cmd{ajPICT\jpincmd*{温泉}} & \ajHotSpring     \\
  \end{booktabular}
\end{center}

\vfil

\subsection{和文合字}

\begin{center}
  \begin{booktabular}{l*{2}{p{13zw}}l}
    \ajMasu            & \cmd{ajMasu}        & \cmd{ajLig\{\jpincmd*{升}\}}   & \ajMasu            \\
    \ajYori            & \cmd{ajYori}        & \cmd{ajLig\{\jpincmd*{より}\}} & \ajYori            \\
    \ajKoto            & \cmd{ajKoto}        & \cmd{ajLig\{\jpincmd*{コト}\}} & \ajKoto            \\
    \ajUta             & \cmd{ajUta}         & --                             & \ajUta             \\
    \ajBlackSesame     & \cmd{ajBlackSesame} & --                             & \ajBlackSesame     \\
    \ajWhiteSesame     & \cmd{ajWhiteSesame} & --                             & \ajWhiteSesame     \\
    \ajJIS             & \cmd{ajJIS}         & \cmd{ajLig\{JIS\}}             & \ajJIS             \\
    {\ibmfamily\ajJAS} & \cmd{ajJAS}         & \cmd{ajLig\{JAS\}}             & {\ibmfamily\ajJAS} \\
  \end{booktabular}
\end{center}

\newpage
\subsection{踊り字}

以下の踊り字は縦書きのみで利用できます。ただし、くの字点は\UCN{3031}、\UCN{3032}のような1字ではなく、\UCN{3033}または\UCN{3034}と\UCN{3035}を組み合わせたものを使用しています。

\begin{center}
  \begin{booktabular}{p{10zw}c}
    \raisebox{0.5zw}{\cmd{ajKunoji}}         & \raisebox{-0.2zw}{\pbox<t>{\ajKunoji}}         \\[0.3zw]
    \raisebox{0.5zw}{\cmd{ajKunojiwithBou}}  & \raisebox{-0.2zw}{\pbox<t>{\ajKunojiwithBou}}  \\[0.3zw]
    \raisebox{0.5zw}{\cmd{ajDKunoji}}        & \raisebox{-0.2zw}{\pbox<t>{\ajDKunoji}}        \\[0.3zw]
    \raisebox{0.5zw}{\cmd{ajDKunojiwithBou}} & \raisebox{-0.2zw}{\pbox<t>{\ajDKunojiwithBou}} \\
  \end{booktabular}
  \hfil
  \begin{booktabular}{p{10zw}c}
    \cmd{ajNinoji}         & \raisebox{-0.2zw}{\pbox<t>{\ajNinoji}}              \\
    \cmd{ajvarNinoji}      & \raisebox{-0.2zw}{\ibmfamily\pbox<t>{\ajvarNinoji}} \\
    \cmd{ajYusuriten}      & \raisebox{-0.2zw}{\ibmfamily\pbox<t>{\ajYusuriten}} \\
  \end{booktabular}
\end{center}

黒ゴマ傍点付きのくの字点(\cmd{aj(D)KunojiwithBou})を利用するには\pkg{plext}が必須です。これは内部で\pkg{plext}から提供される\cmd{bou}を利用して黒ゴマ(\cmd{ajBlackSesame})を与えているためです。\warichu*{エラーの際に\pkg{plext}が必要な旨のメッセージはありませんが、&\cmd{bou}を定義しているパッケージは\pkg{plext}のみです}

これを踏まえると、\cmd{boutenchar}を再定義することで黒ゴマ傍点から白ゴマ傍点に変更することも可能です。

\begin{demosource}*[0.2]
  {\textbackslash}ajKunojiwithBou%
  {\textbackslash}ajDKunojiwithBou\\
  \vfil
  {\textbackslash}renewcommand\{{\textbackslash}boutenchar\}\{{\textbackslash}hbox to 1zw\{{\textbackslash}tate{\textbackslash}hfil{\textbackslash}tiny{\textbackslash}ajWhiteSesame\}\}\\
  \vfil
  {\textbackslash}ajKunojiwithBou%
  {\textbackslash}ajDKunojiwithBou
  \tcblower
  \pbox<t>{\ajKunojiwithBou}
  \hspace{1zw}
  \pbox<t>{\ajDKunojiwithBou}

  \vspace{-0.6zw}
  {\color{gray9}\rule{\linewidth}{0.15zw}}\\[-1zw]
  {\scriptsize (\cmd{boutenchar}の再定義)}\\[-1.3zw]
  {\color{gray9}\rule{\linewidth}{0.15zw}}
  \vspace{-1zw}

  \renewcommand{\boutenchar}{\hbox to 1zw{\tate\hfil\tiny\ajWhiteSesame}}%

  \pbox<t>{\ajKunojiwithBou}
  \hspace{1zw}
  \pbox<t>{\ajDKunojiwithBou}
\end{demosource}

\vfil

\subsection{その他}

\begin{center}
  \begin{booktabular}{p{10zw}>{\ibmfamily}c}
    \cmd{ajSenteMark}    & \ajSenteMark     \\
    \cmd{ajGoteMark}     & \ajGoteMark      \\
    \cmd{ajCheckmark}    & \ajCheckmark     \\
    \cmd{ajVisibleSpace} & \ajVisibleSpace  \\
    \cmd{ajCommandKey}   & \ajCommandKey    \\
    \cmd{ajReturnKey}    & \ajReturnKey     \\
  \end{booktabular}
\end{center}

\newpage
\section{リーダー：\texorpdfstring{\cmd{ajLeader}}{\textbackslash ajLeader}}

区切り線等のリーダーを記述するための\cmd{ajLeader}があります。そのままでは\cmd{hfill}のように文字で行末まで埋めますが、オプション引数で\eghostguarded{\textlrangle{dim}}を指定することでその幅で出力することも出来ます。

\begin{demosource}
  ここから{\textbackslash}ajLeader\{2\}ここまで\\
  ここから{\textbackslash}ajLeader[3zw]\{3\}ここまで
  \tcblower
  \ibmfamily
  ここから
  \ajLeader{2}%
  ここまで\\
  ここから
  \ajLeader[3zw]{3}%
  ここまで
\end{demosource}

\begin{center}
  \begin{booktabular}{>{\ttfamily}c>{\ibmfamily}p{10zw}}
    \thead{引数} & \thead{サンプル} \\
    \midrule
    1  & \makebox[10zw]{\ajLeader[10zw]{1}}    \\
    2  & \makebox[10zw]{\ajLeader[10zw]{2}}    \\
    3  & \makebox[10zw]{\ajLeader[10zw]{3}}    \\
  \end{booktabular}
  \hfil
  \begin{booktabular}{>{\ttfamily}c>{\ibmfamily}p{10zw}}
    \thead{引数} & \thead{サンプル} \\
    \midrule
    4  & \makebox[10zw]{\ajLeader[10zw]{4}}    \\
    5  & \makebox[10zw]{\ajLeader[10zw]{5}}    \\
    6  & \makebox[10zw]{\ajLeader[10zw]{6}}    \\
  \end{booktabular}
\end{center}

6を超える値を与えると、例えば\cmd{ajLeader\{7\}}では{\ibmfamily\ajLeader[4.6zw]{7}}のような記号が連続するようになります。ただし、横組と縦組で記号が異なるようになります。\warichu{本ドキュメントではリーダーとして利用されるのは6までだろうと判断しています}

\section{括弧：\texorpdfstring{\cmd{ajQuote}}{\textbackslash ajQuote}}

\cmd{ajQuote}を利用することで、簡単にその値に応じた括弧を出力できます。

\begin{demosource}
  {\textbackslash}ajQuote\{11\}\{朗報\}
  \tcblower
  \ajQuote{11}{朗報}
\end{demosource}

デフォルトでは以下の表の通りの値と括弧になります。19以上の数値を与えるとエラーとなります。

\newcommand{\ajQuoteS}[1]{#1 & \ibmfamily\ajQuote{#1}{○} & \raisebox{0pt}{\parbox<t>{3zw}{\ibmfamily\ajQuote{#1}{○}}}}
\newcommand{\breakcell}[1]{\setlength{\tabcolsep}{0pt}\hspace{-\tabcolsep}\( \Bigl\{\begin{tabular}{l} #1 \end{tabular} \)}

\begin{center}
  \begin{booktabular}{>{\ttfamily}cccp{10zw}}
    \thead{値} & \thead{横組} & \thead{縦組} & \thead{括弧種} \\
    \midrule
    \ajQuoteS{1}  & \breakcell{\scriptsize 横：シングルクォート\\[-0.55zw]\scriptsize 縦：シングルミニュート} \\[-0.5zw]
    \ajQuoteS{2}  & \breakcell{\scriptsize 横：ダブルクォート\\[-0.55zw]\scriptsize 縦：ダブルミニュート} \\[-0.5zw]
    \ajQuoteS{3}  & 丸括弧 \\[-0.5zw]
    \ajQuoteS{4}  & 亀甲括弧 \\[-0.5zw]
    \ajQuoteS{5}  & 角括弧 \\[-0.5zw]
    \ajQuoteS{6}  & 波括弧 \\[-0.5zw]
    \ajQuoteS{7}  & 山括弧 \\[-0.5zw]
    \ajQuoteS{8}  & 二重山括弧 \\[-0.5zw]
    \ajQuoteS{9}  & かぎ括弧 \\[-0.5zw]
  \end{booktabular}
  \hfil
  \begin{booktabular}{>{\ttfamily}cccp{10zw}}
    \thead{値} & \thead{横組} & \thead{縦組} & \thead{括弧種} \\
    \midrule
    \ajQuoteS{10} & 二重かぎ括弧 \\[-0.5zw]
    \ajQuoteS{11} & 隅付き括弧 \\[-0.5zw]
    \ajQuoteS{12} & ダブルミニュート \\[-0.5zw]
    \ajQuoteS{13} & 二重かぎ括弧 \\[-0.5zw]
    \ajQuoteS{14} & 二重角括弧 \\[-0.5zw]
    \ajQuoteS{15} & 二重亀甲括弧 \\[-0.5zw]
    \ajQuoteS{16} & 二重丸括弧 \\[-0.5zw]
    \ajQuoteS{17} & シングルミニュート \\[-0.5zw]
    \ajQuoteS{18} & 白隅付き括弧 \\[-0.5zw]
  \end{booktabular}
\end{center}

縦組の\cmd{ajQuote\{1\}}や\cmd{ajQuote\{2\}}がミニュートになるのは\pkg{otf}のフォントメトリックとは独立しています。そのため、\cls{jlreq}等の他のフォントメトリックを利用している場合であっても縦組の際にはミニュートになります。

\newpage
\section{漢文：\texorpdfstring{\cmd{ajKunten}}{\textbackslash ajKunten}}

漢文で使用される訓点には、ふつうに入力される文字とは異なる特別な文字があります。これらの組はほとんど見分けが付きませんが別の文字になっています。
また、「\ajKunten{一}\subUCN{3192}」と「\ajKunten{レ}\subUCN{3191}」を繋げた一レ点のようなものもあります。
これらの訓点は\cmd{ajKunten\{\textlrangle{引数}\}}を利用して表現します。

\newcommand{\ajKuntenS}[3]{#2 {\rmfamily\normalsize\subUCN{#1}} & \ajKunten{#2} \subUCN{#3}}

\begin{center}
  \begin{booktabular}{>{\ttfamily\small}C{5zw}C{5zw}}
    \thead{引数} & \thead{訓点} \\
    \midrule
    \ajKuntenS{30EC}{レ}{3191} \\
    \ajKuntenS{4E00}{一}{3192} \\
    \ajKuntenS{4E8C}{二}{3193} \\
    \ajKuntenS{4E09}{三}{3194} \\
    \ajKuntenS{56DB}{四}{3195} \\
    \ajKuntenS{4E0A}{上}{3196} \\
    \ajKuntenS{4E2D}{中}{3197} \\
    \ajKuntenS{4E0B}{下}{3198} \\
  \end{booktabular}
  \hfil
  \begin{booktabular}{>{\ttfamily\small}C{5zw}C{5zw}}
    \thead{引数} & \thead{訓点} \\
    \midrule
    \ajKuntenS{7532}{甲}{3199} \\
    \ajKuntenS{4E59}{乙}{319A} \\
    \ajKuntenS{4E19}{丙}{319B} \\
    \ajKuntenS{4E01}{丁}{319C} \\
    \ajKuntenS{5929}{天}{319D} \\
    \ajKuntenS{5730}{地}{319E} \\
    \ajKuntenS{4EBA}{人}{319F} \\
  \end{booktabular}
  \hfil
  \begin{booktabular}{>{\ttfamily\small}C{5zw}C{5zw}}
    \thead{引数} & \thead{訓点} \\
    \midrule
    一レ & \large\raisebox{-0.2zw}{\ajKunten{一レ}} \\[1zw]
    上レ & \large\raisebox{-0.5zw}{\ajKunten{上レ}} \\[1zw]
    甲レ & \large\raisebox{-0.5zw}{\ajKunten{甲レ}} \\[1zw]
    天レ & \large\raisebox{-0.5zw}{\ajKunten{天レ}} \\
  \end{booktabular}
\end{center}

以上19の訓点の記法には、\cmd{aj\jpincmd*{訓点レ}}や\cmd{aj\jpincmd*{訓点上レ}}のように\cmd{aj\jpincmd*{訓点}\textlrangle{引数}}とするコマンドも提供されています。

これらに加えて、「｜\subUCN{FF5C}」を利用して\cmd{ajKunten\{｜\}}とすることで\ruby{竪点}{たて|てん}（連字符号）「\ajKunten{｜}\subUCN{3190}」も表現できます。
また、竪点と返り点を組み合わせた訓点も表現できます。

\newcommand{\ajKuntenVS}[1]{#1 & \ajKunten{#1}}

\begin{center}
  \begin{booktabular}{>{\ttfamily\small}C{3zw}C{3zw}}
    \thead{引数} & \thead{訓点} \\
    \midrule
    \ajKuntenVS{二｜} \\
    \ajKuntenVS{三｜} \\
    \ajKuntenVS{四｜} \\
    \ajKuntenVS{中｜} \\
    \ajKuntenVS{下｜} \\
  \end{booktabular}
  \hfil
  \begin{booktabular}{>{\ttfamily\small}C{3zw}C{3zw}}
    \thead{引数} & \thead{訓点} \\
    \midrule
    \ajKuntenVS{乙｜} \\
    \ajKuntenVS{丙｜} \\
    \ajKuntenVS{丁｜} \\
    \ajKuntenVS{地｜} \\
    \ajKuntenVS{人｜} \\
  \end{booktabular}
  \hfil
  \begin{booktabular}{>{\ttfamily\small}C{3zw}C{3zw}}
    \thead{引数} & \thead{訓点} \\
    \midrule
    \ajKuntenVS{｜二} \\
    \ajKuntenVS{｜三} \\
    \ajKuntenVS{｜四} \\
    \ajKuntenVS{｜中} \\
    \ajKuntenVS{｜下} \\
  \end{booktabular}
  \hfil
  \begin{booktabular}{>{\ttfamily\small}C{3zw}C{3zw}}
    \thead{引数} & \thead{訓点} \\
    \midrule
    \ajKuntenVS{｜乙} \\
    \ajKuntenVS{｜丙} \\
    \ajKuntenVS{｜丁} \\
    \ajKuntenVS{｜地} \\
    \ajKuntenVS{｜人} \\
  \end{booktabular}
\end{center}

最後に、\pkg{gckanbun}による例を紹介します。このパッケージでは\cmd{gckanbunkaeriten}によって訓点を表すので、次のように別に定義すると簡便になります。

\begin{demosource}*[0.1]
  {\textbackslash}newcommand\{{\textbackslash}Kunten\}[1]\{{\textbackslash}gckanbunkaeriten\{{\textbackslash}ajKunten\{\#1\}\}\}\\
  \vfil
  有{\textbackslash}Kunten\{\jpincmd{下}\}%
  鬻{\textbackslash}Kunten\{\jpincmd{二}\}%
  盾与{\textbackslash}Kunten\{\jpincmd{一レ}\}%
  矛者{\textbackslash}Kunten\{\jpincmd{上}\}%
  \tcblower
  \newcommand{\Kunten}[1]{\gckanbunkaeriten{\ajKunten{#1}}}
  \LARGE
  \pbox<t>{ 有\Kunten{下}鬻\Kunten{二}盾与\Kunten{一レ}矛者\Kunten{上} }
\end{demosource}

\subsection{\texorpdfstring{\filePath{sfkanbun.sty}}{sfkanbun.sty}、\texorpdfstring{\filePath{kunten2e.sty}}{kunten2e.sty}との連携}

漢文に関する{\pLaTeX}向けパッケージとして\filePath{sfkanbun.sty}と\filePath{kunten2e.sty}があります。これらのパッケージでは、レ点や踊り字がフォントに収録されていないことを受けて、（今からするとかなり強引な）代替方法が提供されています。

\pkg{otf}では、以下の2つのコマンドによって漢文における返り点等のコマンド
（\begin{itemize*}[label={}, afterlabel={}, itemjoin=、]
  \item \cmd{kundokusize}
  \item \cmd{kaeriten\{\textlrangle{返り点}\}}
  \item \cmd{kokana\{\textlrangle{仮名}\}\{\textlrangle{返り点}\}}
  \item \cmd{reten}
  \item \cmd{retenform}
  \item \cmd{retenkana\{\textlrangle{レ点付き仮名・仮名}\}}
\end{itemize*}）
のスタイルを{\AdobeJapan}にあるグリフ変更するためのコマンドを提供します。

\begin{itemize}
  \item \cmd{DeclareOriginalKundokuStyle}
  \item \cmd{DeclareAJKundokuStyle}
\end{itemize}

これらのコマンドは以上で示しているパッケージを前提としたものになっているため、2025年現在では有用なコマンドとは言い難い状況となっています。

\vfil

ここでは最後に2つのパッケージに関する情報を以下に示します。\warichu{2つともライセンス周りが明示されていないため、本リポジトリでの再配布は避けています}

\begin{itemize}
  \item {\large\filePath{sfkanbun.sty}}（\ruby{藤田}{ふじ|た}\ruby{眞作}{しん|さく}{\氏}）

  パッケージは以下のページから取得できます。\warichu*{ただし、\filePath{sfkanbun.sty}は内部で&\filePath{jdkintou.sty}が必要になります}

  {\noindent}TeX/LaTeX Applications by Shinsaku Fujita - Fujita Shinsaku TeX/LaTeX Page
  \urltab{\url{http://xymtex.com/fujitas2/texlatex/index.html}}

  ただし、これに関するマニュアルはありません。ソース内にある程度の解説があるため、これを読めば凡そ使えるようになります。
  \vspace{1zw}
  \item {\large\filePath{kunten2e.sty}}（\ruby{金水}{きん|すい}\ruby{敏}{さとし}{\氏}）

  パッケージは以下の大阪大学のページから配布されていましたが、2025年10月現在、アクセスすることは出来ずパッケージを取得できません。\warichu{2022年度末に大阪大学を退職されているので、それに伴ってページが削除された？}
  \urltab[network-off]{\url{http://www.let.osaka-u.ac.jp/~kinsui/tex/top.htm}}

  \filePath{kunten2e.sty}を取得する代替として、Wayback Machineを利用する方法があります。ここから\filePath{kunten2e.sty}ファイルとパッケージの解説PDF「kunten2e.styについて」(\filePath{kunkai.pdf})を取得できます。
  \urltab{\url{https://web.archive.org/web/20220210173434/http://www.let.osaka-u.ac.jp/~kinsui/tex/top.htm}}
  \vspaceLine

  \filePath{kunkai.pdf}はマニュアルに相当しますが、公開されているバージョン2.0の完全なマニュアルとはなっていません。
\end{itemize}

ただし、これらのパッケージはそれなりに古く、仮に今すべてのパッケージを取得できたとしても今の{\LaTeX}で動作する保証がないことに注意してください。

特に、\filePath{sfkanbun.sty}に関しては\href{https://okumuralab.org/tex/mod/forum/discuss.php?d=2661}{{\TeX} Forum \#\textsf{2661}}を参照してください。2018年12月以降の{\LaTeX}で生じる問題を修正する方法と修正済みのファイルを取得できます。

\end{document}
